% Options for packages loaded elsewhere
\PassOptionsToPackage{unicode}{hyperref}
\PassOptionsToPackage{hyphens}{url}
%
\documentclass[
]{article}
\usepackage{amsmath,amssymb}
\usepackage{iftex}
\ifPDFTeX
  \usepackage[T1]{fontenc}
  \usepackage[utf8]{inputenc}
  \usepackage{textcomp} % provide euro and other symbols
\else % if luatex or xetex
  \usepackage{unicode-math} % this also loads fontspec
  \defaultfontfeatures{Scale=MatchLowercase}
  \defaultfontfeatures[\rmfamily]{Ligatures=TeX,Scale=1}
\fi
% \usepackage{lmodern} (commented; using unicode-fix)
% unicode-fix.tex — render raw-Unicode symbols that the Latin Modern *text* font
% lacks, using Computer Modern math/text equivalents. This keeps the established
% Computer Modern look while filling the glyphs xelatex would otherwise drop.
% Included via: pandoc ... --include-in-header=unicode-fix.tex
\usepackage{newunicodechar}
\usepackage{pifont}
% --- Greek letters -> math (Latin Modern Math; consistent with CM body text) ---
\newunicodechar{α}{\ensuremath{\alpha}}
\newunicodechar{β}{\ensuremath{\beta}}
\newunicodechar{γ}{\ensuremath{\gamma}}
\newunicodechar{δ}{\ensuremath{\delta}}
\newunicodechar{ε}{\ensuremath{\varepsilon}}
\newunicodechar{θ}{\ensuremath{\theta}}
\newunicodechar{μ}{\ensuremath{\mu}}
\newunicodechar{ν}{\ensuremath{\nu}}
\newunicodechar{ρ}{\ensuremath{\rho}}
\newunicodechar{σ}{\ensuremath{\sigma}}
\newunicodechar{τ}{\ensuremath{\tau}}
\newunicodechar{φ}{\ensuremath{\varphi}}
\newunicodechar{χ}{\ensuremath{\chi}}
\newunicodechar{ψ}{\ensuremath{\psi}}
\newunicodechar{η}{\ensuremath{\eta}}
\newunicodechar{κ}{\ensuremath{\kappa}}
\newunicodechar{λ}{\ensuremath{\lambda}}
\newunicodechar{π}{\ensuremath{\pi}}
\newunicodechar{ω}{\ensuremath{\omega}}
\newunicodechar{Ḣ}{\ensuremath{\dot H}}
\newunicodechar{≪}{\ensuremath{\ll}}
\newunicodechar{≫}{\ensuremath{\gg}}
\newunicodechar{⟺}{\ensuremath{\Longleftrightarrow}}
\newunicodechar{⟹}{\ensuremath{\Longrightarrow}}
\newunicodechar{⇔}{\ensuremath{\Leftrightarrow}}
\newunicodechar{⟨}{\ensuremath{\langle}}
\newunicodechar{⟩}{\ensuremath{\rangle}}
\newunicodechar{⁰}{\textsuperscript{0}}
\newunicodechar{¹}{\textsuperscript{1}}
\newunicodechar{²}{\textsuperscript{2}}
\newunicodechar{³}{\textsuperscript{3}}
\newunicodechar{⁴}{\textsuperscript{4}}
\newunicodechar{⁵}{\textsuperscript{5}}
\newunicodechar{⁶}{\textsuperscript{6}}
\newunicodechar{⁷}{\textsuperscript{7}}
\newunicodechar{⁸}{\textsuperscript{8}}
\newunicodechar{⁹}{\textsuperscript{9}}
% --- Relations / operators -> math ---
\newunicodechar{↔}{\ensuremath{\leftrightarrow}}
\newunicodechar{⇐}{\ensuremath{\Leftarrow}}
\newunicodechar{⇒}{\ensuremath{\Rightarrow}}
\newunicodechar{∂}{\ensuremath{\partial}}
\newunicodechar{∅}{\ensuremath{\varnothing}}
\newunicodechar{∈}{\ensuremath{\in}}
\newunicodechar{∼}{\ensuremath{\sim}}
\newunicodechar{≈}{\ensuremath{\approx}}
\newunicodechar{≠}{\ensuremath{\neq}}
\newunicodechar{≤}{\ensuremath{\leq}}
\newunicodechar{≥}{\ensuremath{\geq}}
\newunicodechar{⊂}{\ensuremath{\subset}}
\newunicodechar{⊃}{\ensuremath{\supset}}
% --- Special ---
\newunicodechar{ℏ}{\ensuremath{\hbar}}
\newunicodechar{ℤ}{\ensuremath{\mathbb{Z}}}
% --- Superscript modifier letters (footnote-style markers) -> text superscript ---
\newunicodechar{ʰ}{\textsuperscript{h}}
\newunicodechar{ʲ}{\textsuperscript{j}}
\newunicodechar{ᵃ}{\textsuperscript{a}}
\newunicodechar{ᵇ}{\textsuperscript{b}}
\newunicodechar{ᵈ}{\textsuperscript{d}}
\newunicodechar{ᵉ}{\textsuperscript{e}}
\newunicodechar{ᵍ}{\textsuperscript{g}}
\newunicodechar{ᶠ}{\textsuperscript{f}}
\newunicodechar{⁻}{\textsuperscript{\ensuremath{-}}}
% --- Subscript digits -> text subscript ---
\newunicodechar{₀}{\textsubscript{0}}
\newunicodechar{₁}{\textsubscript{1}}
\newunicodechar{₂}{\textsubscript{2}}
\newunicodechar{₃}{\textsubscript{3}}
\newunicodechar{₄}{\textsubscript{4}}
\newunicodechar{₈}{\textsubscript{8}}
% --- Checkmarks ---
\newunicodechar{✓}{\ding{51}}
\newunicodechar{✗}{\ding{55}}

\ifPDFTeX\else
  % xetex/luatex font selection
\fi
% Use upquote if available, for straight quotes in verbatim environments
\IfFileExists{upquote.sty}{\usepackage{upquote}}{}
\IfFileExists{microtype.sty}{% use microtype if available
  \usepackage[]{microtype}
  \UseMicrotypeSet[protrusion]{basicmath} % disable protrusion for tt fonts
}{}
\makeatletter
\@ifundefined{KOMAClassName}{% if non-KOMA class
  \IfFileExists{parskip.sty}{%
    \usepackage{parskip}
  }{% else
    \setlength{\parindent}{0pt}
    \setlength{\parskip}{6pt plus 2pt minus 1pt}}
}{% if KOMA class
  \KOMAoptions{parskip=half}}
\makeatother
\usepackage{xcolor}
\setlength{\emergencystretch}{3em} % prevent overfull lines
\providecommand{\tightlist}{%
  \setlength{\itemsep}{0pt}\setlength{\parskip}{0pt}}
\setcounter{secnumdepth}{-\maxdimen} % remove section numbering
\ifLuaTeX
  \usepackage{selnolig}  % disable illegal ligatures
\fi
\IfFileExists{bookmark.sty}{\usepackage{bookmark}}{\usepackage{hyperref}}
\IfFileExists{xurl.sty}{\usepackage{xurl}}{} % add URL line breaks if available
\urlstyle{same}
\hypersetup{
  hidelinks,
  pdfcreator={LaTeX via pandoc}}

\author{}
\date{}

\begin{document}

\hypertarget{the-substratum-construction-reconstruction-the-substratum-gauge-group-and-the-qm-gr-synthesis}{%
\section{The Substratum Construction: Reconstruction, the Substratum
Gauge Group, and the QM-GR
Synthesis}\label{the-substratum-construction-reconstruction-the-substratum-gauge-group-and-the-qm-gr-synthesis}}

\textbf{Author:} Alex Maybaum\\
\textbf{Date:} April 2026\\
\textbf{Status:} DRAFT PRE-PRINT\\
\textbf{Classification:} Theoretical Physics / Foundations

\begin{center}\rule{0.5\linewidth}{0.5pt}\end{center}

\hypertarget{abstract}{%
\subsection{Abstract}\label{abstract}}

Quantum mechanics and general relativity are conventionally treated as
two fundamental theories awaiting unification. This paper establishes
that they are instead two projections of a single finite deterministic
construction --- joined by a third projection that produces the arrow of
time --- with the apparent QM-GR incompatibility arising as a category
error rather than a technical problem. The construction is the bijection
\((S, \varphi)\) on a finite configuration space whose visible-sector
projection is the Standard Model {[}SM{]} and whose
boundary-thermodynamic projection is general relativity {[}GR{]}, as
developed in companion papers. The substratum-level results that make
these emergences a single construction rather than two independent
applications of the same formalism are: (i) the \emph{reconstruction
theorem} (Theorem 23), which shows that observed physics --- quantum
mechanics with Bell violations, finite boundary entropy, and spatial
isotropy --- together with the framework's structural assumptions
(finiteness, determinism, bounded coupling, center independence,
linearity, background independence) uniquely determines the equivalence
class \([(S, \varphi)]/\mathcal{G}_{\text{sub}}\) at the lattice level,
with the Standard Model gauge group
\(\mathrm{SU}(3) \times \mathrm{SU}(2) \times \mathrm{U}(1)\) and
\(\bar\theta = 0\) as forced retrodictions, and the anomaly-free
hypercharge assignment forced once the observed family pattern is given
(fermion embedding closed via the link-carrier construction, with the
absolute generation count locked at exactly three by coupling-degree
minimality); and (ii) the \emph{substratum gauge group}
\(\mathcal{G}_{\text{sub}}\) (Theorem 24), with four explicit families
of generators proved to exhaust all observables-preserving
transformations via Stinespring uniqueness. The Standard Model gauge
group is the visible-sector shadow of \(\mathcal{G}_{\text{sub}}\) ---
the image of the substratum's symmetry under the trace-out --- and the
three-level gauge hierarchy (substratum, emergent QFT, emergent
Hamiltonian) provides the structural account of why the Standard Model
has its specific gauge structure. The synthesis claim is that the QM-GR
incompatibility, as conventionally posed, asks how to merge two
descriptions that share no common framework; the present construction
provides that common framework, and the apparent incompatibility is the
artifact of treating two different projections of one object as if they
were two different theories. The result is not a unification of quantum
mechanics and general relativity in the traditional sense --- that
program remains as ill-posed as before --- but a dissolution of the
question.

\begin{center}\rule{0.5\linewidth}{0.5pt}\end{center}

\hypertarget{introduction}{%
\subsection{1. Introduction}\label{introduction}}

The standard physical picture of the world contains two fundamental
theories. Quantum mechanics describes the microscopic regime in terms of
unitary evolution on a Hilbert space, superposition, and Born-rule
probabilities. General relativity describes the gravitational regime in
terms of a deterministic, classical, dynamical metric obeying Einstein's
equations. The two are mathematically and conceptually incompatible.
Quantum mechanics has no preferred temporal foliation and no
observer-independent state; general relativity has both. Quantum
mechanics treats measurement as a primitive operation; general
relativity treats it as a process within the dynamics. Quantum mechanics
is linear; general relativity is not. Every attempt to unify them --- to
quantize gravity, to geometrize quantum mechanics, to merge the two into
a deeper underlying theory --- has produced either inconsistencies,
ambiguities, or empirically unfalsifiable conjectures. The unification
problem has resisted solution for nearly a century, and the most active
current programs (loop quantum gravity, string theory, causal dynamical
triangulations, asymptotic safety) have produced neither a quantitative
empirical confirmation nor a clean conceptual resolution.

This paper argues that the QM-GR incompatibility is not a technical
problem to be solved by finding the right mathematical framework. It is
a category error: an attempt to merge two descriptions that turn out to
be projections of the same underlying object viewed from different
perspectives, rather than two competing accounts of the same regime. The
framework developed in this four-paper synthesis to which this paper
belongs --- {[}Main{]}, {[}SM{]}, {[}GR{]}, and the present paper ---
provides the underlying object and the projection map. The underlying
object is a finite deterministic bijection \((S, \varphi)\) on a
configuration space partitioned into a visible sector accessible to an
embedded observer and a hidden sector beyond the observer's causal
reach. The projection from the underlying object to quantum mechanics is
the trace-out over the hidden sector, formalized in {[}Main{]} as the
embedded-observation theorem: under three structural conditions (C1:
non-trivial coupling, C2: slow bath, C3: high-capacity hidden sector),
the embedded observer's description is necessarily quantum mechanics,
with the wave function, Born rule, and unitary evolution arising as
derived rather than fundamental constructs. The projection from the same
underlying object to general relativity is the thermodynamic limit at
the partition boundary, formalized in {[}GR{]} as the
cosmological-horizon application: applied to the cosmological horizon as
a causal partition, the framework determines
\(\hbar = c^3 (2 l_p)^2 / (4G)\), the Bekenstein-Hawking entropy with
the \(1/4\) coefficient, and the dissolution of the cosmological
constant problem. The Standard Model emerges from the same bijection on
a cubic lattice with the wave equation as substratum dynamics, with the
gauge group
\(\mathrm{SU}(3) \times \mathrm{SU}(2) \times \mathrm{U}(1)\), three
chiral generations, the Higgs as composite, and \(\bar\theta = 0\) as
forced consequences {[}SM{]}. The framework's hierarchical structural
realism and its relationship to other unification programs are developed
in the fifth core paper {[}Structure{]}.

The present paper develops the substratum-level results that make these
three emergences --- emergent QM in {[}Main{]}, the Standard Model
derivation in {[}SM{]}, and the gravitational sector in {[}GR{]} --- a
single construction rather than three separate applications of the same
formalism. The two technical results that do this work are the
\emph{reconstruction theorem} (§3) and the \emph{substratum gauge group}
(§4). The reconstruction theorem shows that the map from
\((S, \varphi)\) to observed physics is invertible within the
framework's structural class: observed quantum mechanics with Bell
violations, together with finite boundary entropy and spatial isotropy
and the structural assumptions of the framework (finiteness,
determinism, bounded coupling, center independence, linearity,
background independence), uniquely determine the equivalence class of
bijections on a lattice, with the Standard Model structure as a forced
retrodiction. The substratum gauge group identifies the kernel of this
inverse map --- four families of transformations of \((S, \varphi)\)
that exhaust all observables-preserving operations --- and shows that
the Standard Model gauge group is its shadow under the trace-out. The
combination of these two results turns the framework's three derivations
into a single object: the SM gauge group is not chosen from a landscape,
the gravitational thermodynamics are not appended to a quantum
description, and the emergent QM is not postulated as a starting point.
All three emerge from the same \((S, \varphi)\), and the substratum
gauge group is the symmetry that makes the relationship between them
exact rather than approximate.

The synthesis claim that follows from these results --- developed in §5
--- is that quantum mechanics and general relativity are not two
theories awaiting unification but two projections of the same
construction, viewed at different levels of description. Quantum
mechanics lives at the level of the visible-sector trace-out: the
embedded observer sees the bijection through the hole of finite causal
access, and the resulting compressed description is unitary QM. General
relativity lives at the level of the boundary thermodynamics: the same
causal partition that produces QM at the visible-sector level produces
classical horizon thermodynamics at the boundary level, and Jacobson's
thermodynamic argument promotes this to Einstein's equations. The two
descriptions are not in tension because they refer to different objects
in the construction --- not different \emph{aspects of the same regime},
which would invite unification, but different \emph{projections of the
same substratum}, which makes unification a category error. The
substratum gauge group makes this precise: at the level of
\((S, \varphi)\), there is no quantum and no classical, no metric and no
wave function; there is only the bijection and its symmetries. The
trace-out at the visible-sector level produces one set of derived
structures (the wave function, the SM gauge group, particle content),
and the thermodynamic limit at the boundary level produces another (the
metric, \(\hbar\), the BH entropy, the dynamical dark energy), and the
framework provides the construction that makes both of these
descriptions exact projections of the same object. The same construction
produces a third projection (§5.4) that resolves the arrow-of-time
question by descending the observer's emergent arrows --- cosmological,
thermodynamic, memory, measurement --- from the horizon-growth clock and
the C1--C3 observer-selection theorem of {[}Main, §4.6{]}, answering the
Ellis objection to bijective substrates without adding a structural
commitment.

The paper is organized as follows. §1.1 states the framework's primary
metaphysical commitment and its relation to explanatory and empirical
claims. §2 develops the physical interpretation of \((S, \varphi)\) as a
finite lossless memory with partial read-write access, and addresses the
substrate objection (what is the memory made of?) by identifying it as
gauge in the precise technical sense established by the substratum gauge
group of §4. §3 states and proves the reconstruction theorem (Theorem
23). §4 states and proves the substratum gauge group result (Theorem 24)
and develops the three-level gauge hierarchy. §5 makes the synthesis
claim explicit, citing {[}Main{]}, {[}SM{]}, and {[}GR{]} for the
empirical content. §6 discusses structural realism, the ontological
hierarchy, and the dissolution of the measurement problem. §7 concludes.
The paper is short by design: the technical content fits in two sections
(§§3--4) and the rest is the structural and interpretive context that
shows why those two sections matter.

\hypertarget{the-metaphysical-commitment}{%
\subsubsection{1.1 The metaphysical
commitment}\label{the-metaphysical-commitment}}

The framework's primary metaphysical claim is \emph{ontic structural
realism applied to the equivalence class}
\([(S, \varphi)] / \mathcal{G}_{\text{sub}}\). What is real at this
level is the equivalence class, not any specific substratum
representative. This is the strongest claim the framework's own theorems
license at the substratum-class level, and it is forced by them: the
reconstruction theorem (§3) shows that the equivalence class is uniquely
determined by the observations together with the structural assumptions,
and the substratum gauge group (§4) shows that different representatives
of the class are observably indistinguishable. The class is what the
data pin down within OI's structural class; individual substrates are
gauge representations of it. A stronger claim --- realism about a
specific bijection --- commits to features the framework's own gauge
group identifies as unobservable. A weaker claim --- pure explanatory
economy without ontic commitment --- undersells what the uniqueness
theorem actually establishes.

The structural-realism commitment operates at multiple levels. Within
OI's specific universality class, the framework's theorems license
realism about \([(S, \varphi)] / \mathcal{G}_{\text{sub}}\). Across
structural classes --- at the level of partial-trace observational
features that any embedded-observer system must respect --- broader
structural realism applies, with the framework's universality-class
equivalence ({[}Structure §9{]}) being the appropriate gauge
equivalence. The full multi-level structure is articulated in
{[}Structure{]}, which canonicalizes both the observation hierarchy and
the gauge hierarchy and their orthogonality, and develops the
cross-framework analysis through universality classes of embedded
observers. The substratum-class commitment of this paper is one level of
a broader hierarchical structural realism: realism at multiple
combinations of observation-hierarchy depth and gauge-hierarchy breadth,
with the framework's empirical content concentrated at the deepest
observation level (Level D in {[}Structure §2{]}) and the most
class-specific gauge level (Level G3 in {[}Structure §2{]}). The
articulation here focuses on this concentrated level because that is
where the reconstruction theorem and substratum gauge group operate; the
broader hierarchical content is developed in {[}Structure{]}.

Two positions often proposed as alternatives to structural realism are,
in this framework, consequences of it. \emph{Explanatory economy}: if
the equivalence class is real, the emergence of quantum mechanics in
{[}Main §3{]} is a theorem about representations of a real structure,
not a coincidence requiring additional interpretation, and the
framework's reduction of postulates follows from the uniqueness rather
than standing as an independent virtue. \emph{Empirical predictions}: if
the class's structure is real, its structural features --- the dimension
\(d = 3\), the \(K = 2d = 6\) internal components, the cubic
decomposition \((3, 2, 1)\), the wave equation, the cosmological-horizon
partition --- imply specific numbers at the observable level, giving the
twenty-two quantitative predictions of {[}SM §7{]} and the gravitational
predictions of {[}GR §7{]}. The three framings (realist, explanatory,
empirical) answer three different questions --- what is real, what is
gained over standard QM, what is testable beyond QM --- rather than
providing competing answers to the same question. The rest of this paper
develops the structural and explanatory consequences; {[}SM{]} and
{[}GR{]} develop the empirical consequences.

\begin{center}\rule{0.5\linewidth}{0.5pt}\end{center}

\hypertarget{the-physical-interpretation-of-s-ux3c6}{%
\subsection{2. The Physical Interpretation of (S,
φ)}\label{the-physical-interpretation-of-s-ux3c6}}

\textbf{Storage and memory.} S is the set of all distinguishable states:
\emph{finite capacity}. φ is a bijection: \emph{perfect memory} ---
information is never created or destroyed. Together, (S, φ) is a finite
lossless memory. The partition V defines the observer's \emph{access}
--- both read and write. The observer reads the visible sector
(measuring x) and writes to the hidden sector (each visible-sector
operation imprints correlations on H through the coupling
\(H_{\text{int}}\)). The slow bath (C2) preserves those writes;
subsequent reads retrieve them --- producing the information backflow
that constitutes P-indivisibility. Quantum mechanics is the statistics
of this read-write cycle: not passive observation of a static memory,
but the self-consistent description that emerges when a finite subsystem
both reads from and writes to a lossless register it cannot fully
access. The Born rule is the equilibrium of the cycle, not of passive
reading. Interference is write-then-read: information deposited in the
hidden sector during one transition is retrieved on a later transition
and partially cancels or reinforces the transition probabilities.

The \(10^{122}\) CC discrepancy is the compression ratio between total
storage and readable storage. The dark sector is the gravitational
effect of the unreadable storage. The Bekenstein-Hawking entropy is the
storage capacity of the partition boundary. The action scale ℏ is the
conversion factor between storage geometry and read-write statistics.

\textbf{Remark (the domain of the substratum dynamics).} The wave
equation of {[}SM §4.1{]} --- the unique second-order reversible
nearest-neighbor dynamics compatible with center independence, isotropy,
and linearity --- is defined on the full substratum \((S, \varphi)\),
not on the visible sector alone. Both visible and hidden sectors run the
same wave equation; the partition \(V\) imposes an observer-access
structure on this common dynamics but does not change it. This
convention makes three pieces of the framework consistent: (i) the
trace-out of {[}Main{]} proceeds by marginalizing over the hidden
sector's substratum degrees of freedom, which requires those degrees of
freedom to have a well-defined dynamics; (ii) the nested trace-out of
{[}GR §8.4{]} extends the construction to a deep hidden sector by
running the same wave equation at all levels; (iii) the link-carrier
construction of {[}SM §4.7.1.2{]} places matter on links that span
visible and hidden sublattices --- which makes sense only because both
sublattices are governed by the same dynamics. The full-substratum wave
equation is therefore the common dynamical input to all three companion
papers, with each applying a different projection map to extract
observer-level physics.

\textbf{The substrate objection.} ``What is the memory made of?'' (S, φ)
is a \emph{complete description} of reality --- it determines all
observables. Space, time, matter, and energy are derived from (S, φ), so
they cannot appear in its definition without circularity. Whether (S, φ)
\emph{is} reality or \emph{describes} reality is provably undecidable by
any measurement (reconstruction theorem below).

\textbf{Relation to computation.} A Turing machine has a tape (storage),
a head (partial read/write access), and a transition function (update
rule). The correspondence is suggestive: S is the tape, V is the head's
access window, φ is the transition function. But three differences are
physically significant. First, a Turing machine's tape is potentially
infinite; S is finite --- finiteness is essential for the recurrence
proof of P-indivisibility and QM, though the effective finiteness result
({[}GR, §8.4{]}) shows the deep hidden sector may be infinite without
affecting the observable physics. Second, a Turing machine is generally
irreversible (it can erase, overwrite, and halt); φ is a bijection ---
nothing is erased, nothing is created, there is no halting. Third, a
Turing machine computes an \emph{extrinsic function} (input → output for
an external user); (S, φ) computes no extrinsic function --- it is a
closed permutation that cycles through states and returns. The
appearance of dynamics, probability, particles, and forces is entirely
the observer's perspective --- what the permutation looks like through
the partial window V.

\textbf{Remark (the Turing connection under effective finiteness).} The
three differences soften under the effective finiteness result. With the
deep hidden sector potentially infinite, the first difference is between
the formal definition (S finite) and the physical requirement (only
\(\mathcal{C}_V \times \mathcal{C}_B\) need be finite). The second
difference is a specialization, not an opposition: reversible Turing
machines (Bennett, Fredkin-Toffoli) are a well-studied subclass. The
third difference --- extrinsic vs.~intrinsic --- is the one that does
physical work. The mapping is then structural: V is the head, H is the
tape, φ is a reversible transition function, and C1--C3 characterize the
architecture. The framework extends Turing's question: instead of asking
what a machine can compute for an external observer, it asks what
computation looks like to a component of the machine --- and proves the
answer is quantum mechanics.

\textbf{The arrow of time.} The substratum has no arrow of time --- φ
and φ⁻¹ are equally valid. Observers satisfying C1--C3 are nonetheless
structurally restricted to the non-equilibrium phase of
\((S, \varphi)\): the theorem in {[}Main, §4.6{]} shows that
\(\tau_B(V)\) is bounded by the local mixing time on equilibrium
microstates, so C2 fails there. This replaces the low-entropy ``past
hypothesis'' with a structural observer-selection rule and excludes
Boltzmann brains structurally rather than anthropically. Entropy
increase is then emergent as a coarse-grained description of the
non-equilibrium dynamics observers necessarily see.

\textbf{The incompleteness family.} The framework's central result
belongs to a family of impossibility-with-structure theorems. Gödel: a
formal system cannot prove all truths about itself --- the unprovable
truths have rigid structure. Turing: a computer cannot decide all
questions about its own behavior --- the undecidable problems have rigid
structure. OI: an observer embedded in a deterministic system cannot
access the complete state --- the emergent description has rigid
structure (quantum mechanics). The common structure is
\emph{self-reference under finite resources}.

The precise OI analog of the halting problem is: \emph{can the observer
determine the hidden-sector state \(h\)?} The question is well-posed ---
\(h\) has a definite value at every moment, because \((S, \varphi)\) is
deterministic. Different \(h\) values produce different physical
outcomes. But the observer provably cannot determine \(h\): multiple
hidden states are compatible with any visible-sector history, and
transition probabilities \(T_{ij}(t)\) are averages over \(h\). The
structural consequence of this inaccessibility is quantum mechanics ---
just as the structural consequence of the halting problem's
undecidability is computability theory. The framework also identifies a
second class of inaccessible quantities --- the alphabet size \(q\)
({[}SM, §2.7{]}) and the deep-sector cardinality \(|\mathcal{C}_D|\)
({[}GR, §3.2{]}) --- but these are \emph{gauge}, not undecidable:
different values produce identical observables, so the question itself
is physically empty. The hidden state \(h\) is undecidable (real answer,
provably inaccessible). The cardinality \(|\mathcal{C}_D|\) is gauge (no
answer to find).

\textbf{Mathematics and physics.} The trace-out performs a
Jordan-Chevalley projection ({[}SM, Appendix A{]}): it extracts the
semisimple part of the dynamics and erases the nilpotent monodromy.
Physics is the semisimple shadow of mathematics --- the diagonalizable
spectral data, projected by the trace-out and organized by the gauge
group's representation structure. The reconstruction theorem (below)
proves that the mathematical description and the physics determine each
other uniquely up to gauge.

\begin{center}\rule{0.5\linewidth}{0.5pt}\end{center}

\hypertarget{the-reconstruction-theorem}{%
\subsection{3. The Reconstruction
Theorem}\label{the-reconstruction-theorem}}

The forward direction --- from \((S, \varphi)\) to observed physics ---
is established by {[}SM §§3.1, 4--5{]} and the companion paper
{[}Main{]}. The converse question is whether the observed physics
uniquely determines \((S, \varphi)\). The reconstruction proceeds in
three stages, each taking specified inputs and producing specified
outputs with an explicit uniqueness claim. The composite theorem
(Theorem 23) then aggregates the stages.

\hypertarget{inputs-to-the-reconstruction}{%
\subsubsection{3.1 Inputs to the
reconstruction}\label{inputs-to-the-reconstruction}}

The reconstruction takes two kinds of inputs: empirical observations and
structural assumptions. Being explicit about both is essential to the
uniqueness claim.

\textbf{Empirical inputs} (facts about the observed universe):

(E1) \textbf{Unitary quantum mechanics.} The observed physics is quantum
mechanical: states are vectors in a complex Hilbert space, time
evolution is unitary, observables are self-adjoint operators,
measurement outcomes follow the Born rule.

(E2) \textbf{Bell violations.} The observed correlations violate Bell
inequalities, ruling out local hidden-variable theories with
factorizable distributions.

(E3) \textbf{Finite boundary entropy.} The entropy of any bounded region
is finite and scales as the area of its boundary. This is supported by
holographic bounds {[}1, 2{]}; the cosmological horizon has finite area
so the bound applies.

(E4) \textbf{Spatial isotropy.} Observed physics is rotationally
invariant; no spatial direction is preferred.

(E5) \textbf{Propagating gravity.} Gravitational degrees of freedom
propagate (gravitational waves are observed). This requires
\(d \geq 3\), since the Weyl tensor vanishes for \(d \leq 2\). (Used in
Stage 2(a) below; one of the three independent forward filters that
select \(d = 3\).)

(E6) \textbf{Stable matter.} Atoms, planets, and bound states exist on
observed timescales. This requires \(d \leq 3\), since the Coulomb
potential in \(d\) dimensions gives unstable atoms for \(d \geq 4\)
(Ehrenfest {[}10{]}) and gravitational orbits are unstable for
\(d \geq 4\) (Tangherlini {[}11{]}). Without stable matter, no embedded
observers exist to define the partition. (Used in Stage 2(a); second
forward filter for \(d = 3\).)

(E7) \textbf{Boundary-entropy concordance.} The observed cosmological
structure satisfies \(\rho_s / \rho_{\text{crit}} \approx 1\) (spatial
flatness near critical density). In \(d\) spatial dimensions, the
general-\(d\) form of this ratio is
\(\rho_s / \rho_{\text{crit}} = 2/(d-1)\), which equals unity only for
\(d = 3\). (Used in Stage 2(a); third forward filter for \(d = 3\). See
{[}SM §3.2{]} and {[}GR §7.2{]} for the derivation.)

\textbf{Structural assumptions} (restrictions on the class of candidate
substrates \((S, \varphi)\)):

(A1) \textbf{Finiteness.} The configuration space \(S\) is finite.
(Follows from E3 + the holographic bound interpreted as Hilbert-space
dimension cutoff \(\dim \mathcal{H} \leq e^{A/4}\).)

(A2) \textbf{Determinism.} \(\varphi: S \to S\) is a bijection
(deterministic, reversible dynamics).

(A3) \textbf{Bounded coupling degree.} Each site is coupled to a bounded
number of neighbors in \(\varphi\)'s action. (Required for locality and
the emergence of a coupling graph with well-defined dimension.)

(A4) \textbf{Center independence.} The dynamics \(\varphi\) does not
depend on a choice of ``center'' site; equivalently, \(\varphi\)
commutes with lattice translations up to gauge. (Required to derive the
wave equation in Stage 2.)

(A5) \textbf{Linearity.} The wave equation for \(\varphi\) is linear.
(Required to obtain the specific gauge structure in Stage 2; nonlinear
alternatives are not ruled out but would require a separate derivation.)

(A6) \textbf{Background independence.} The dynamics is invariant under
spatially-varying internal-index transformations that preserve the
cubic-symmetric coupling matrix pointwise. The promotion of the
resulting global commutant symmetry to local gauge invariance is then a
derivation step ({[}SM §3.1{]}) --- not part of the assumption itself.

The theorem's uniqueness claim holds under E1--E7 \textbf{and} A1--A6
jointly; removing any of A3--A6 either requires new derivations or
weakens the uniqueness. The set A1--A6 is sufficient for the
reconstruction but is not proved to be minimal: A4 (center independence)
and A6 (background independence) overlap in physical content (A6
promotes the symmetry that A4 constrains), and A5 (linearity) is partly
derived from the other assumptions via the dynamics-selection argument
of {[}SM §4.1{]}. A tighter axiomatization may be possible but is not
pursued here; the reconstruction's validity depends on the sufficiency
of A1--A6 and E1--E7, not on their independence. The empirical inputs
E5, E6, E7 (propagating gravity, stable matter, boundary-entropy
concordance) provide explicit forward filters for \(d = 3\) --- see
Stage 2(a) below.

\textbf{Critical dependencies.} Theorem 23's proof relies on the
following prior theorems, whose individual correctness is assumed:

\begin{itemize}
\tightlist
\item
  {[}Main §2.3{]} P-indivisibility of the marginalized stochastic
  process (Bell violations → P-indivisibility)
\item
  {[}Main §3.2{]} Forward Stinespring dilation (bijection → CPTP
  channel) and reverse-direction lemma (stochastic process → bijection)
\item
  {[}Main §3.4{]} Characterization theorem (Bell violation ⇒ C1, QM ⇒ C2
  + C3, conditional on ETH for C2)
\item
  {[}SM §3.2{]} Coupling-graph dimension → d = 3
\item
  {[}SM §4.1{]} Center independence + isotropy + linearity → wave
  equation
\item
  {[}SM Theorems 5--15{]} Gauge group, generations, hypercharges
  derivation chain
\item
  {[}SM Theorems 17--21{]} T-invariance → \(\bar\theta = 0\)
\item
  {[}GR §3{]} Gap equation from thermal self-consistency →
  \(\hbar = c^3 \epsilon^2 / (4G)\)
\end{itemize}

Theorem 23's confidence is bounded above by the minimum confidence in
these dependencies.

\hypertarget{stage-1-observed-qm-deterministic-embedding-with-c1c3}{%
\subsubsection{3.2 Stage 1: Observed QM → deterministic embedding with
C1--C3}\label{stage-1-observed-qm-deterministic-embedding-with-c1c3}}

\textbf{Inputs:} E1 (unitary QM), E2 (Bell violations), E3 (finite
boundary entropy), A1 (finiteness), A2 (determinism).

\textbf{Output:} There exists a triple \((S, \varphi, V)\) with \(S\)
finite, \(\varphi\) a bijection on \(S\), and \(V \subset S\) a
distinguished subset such that: - The observed quantum dynamics is the
reduced description obtained by tracing out \(S \setminus V\) from the
deterministic evolution under \(\varphi\). - The triple satisfies C1
(non-trivial coupling), C2 (slow-bath memory), and C3 (high
hidden-sector capacity).

\textbf{Conditional structure.} C1 and C3 are derived unconditionally
from the inputs. C2 is conditional on the eigenstate thermalization
hypothesis (ETH) in the hidden sector --- without ETH, the necessity
direction (P-indivisibility on accessible timescales forces
\(\tau_S \ll \tau_B\)) cannot be established, as discussed in {[}Main
§3.4{]}'s ``Remark (Status of ETH).'' ETH is a well-supported conjecture
for generic chaotic many-body Hamiltonians and is the standard
assumption for the cosmological-horizon complement; relaxations or
failures of ETH would require a separate argument. The framework's
gauge-group derivation through Stage 2 inherits this conditional.

\textbf{Derivation.} The observed quantum dynamics provides, for each
choice of computational basis in the visible sector, a family of
transition probabilities \(T_{ij}(t)\) between basis states --- the
Born-rule matrix elements of the unitary evolution. These form a
stochastic process on the finite visible configuration space (finiteness
of \(\mathcal{C}_V\) follows from E3, the boundary entropy bound). Bell
violations (E2) are incompatible with a factorizable hidden-variable
distribution and hence imply the stochastic process is P-indivisible: at
intermediate times no valid stochastic matrix \(\Lambda\) satisfies
\(T(t_2, 0) = \Lambda \, T(t_1, 0)\) ({[}Main §2.3{]}). By the
reverse-direction lemma at the stochastic-process level ({[}Main
§3.2{]}), any such P-indivisible process on a finite configuration space
admits a realization as marginalization of a deterministic bijection
\(\varphi\) on an enlarged finite space
\(\mathcal{C}_V \times \mathcal{C}_H\) with uniform prior on
\(\mathcal{C}_H\). This delivers the triple \((S, \varphi, V)\) with
\(S = \mathcal{C}_V \times \mathcal{C}_H\) and \(V = \mathcal{C}_V\).
Non-permutation \(T(t)\) (forced by Bell violations) gives C1. The
characterization theorem {[}Main §3.4{]} establishes that the resulting
dynamics, conditional on ETH in the hidden sector, necessarily satisfies
C2 (slow-bath memory is required for the non-Markovian returns that
produce quantum interference) and C3 (high hidden-sector capacity is
required for the data-processing bound to accommodate observed
information backflow).

\textbf{Uniqueness at Stage 1.} The triple \((S, \varphi, V)\) is
uniquely determined by the observations up to the equivalences: -
Stinespring ambiguity: different dilations producing the same channel
differ by unitary rotations on the hidden sector, which are absorbed
into \(\mathcal{G}_{\text{sub}}\) (Theorem 24). - Hidden-sector basis
choice: relabeling of hidden states is gauge (generator (i) of
\(\mathcal{G}_{\text{sub}}\)). - Deep-sector enlargement: additional
hidden degrees of freedom with \(\tau_B^D \gg \tau_S\) leave
observations unchanged (generator (iii) of
\(\mathcal{G}_{\text{sub}}\)).

\textbf{Status:} Theorem.

\hypertarget{stage-2-embedding-isotropy-specific-lattice-gauge-structure}{%
\subsubsection{3.3 Stage 2: Embedding + isotropy → specific lattice +
gauge
structure}\label{stage-2-embedding-isotropy-specific-lattice-gauge-structure}}

\textbf{Inputs:} Stage 1 output + E4 (spatial isotropy) + E5
(propagating gravity) + E6 (stable matter) + E7 (boundary-entropy
concordance) + A3 (bounded coupling) + A4 (center independence) + A5
(linearity) + A6 (background independence) + max-lattice-speed condition
(for the wave-equation coupling constant; see step (b) below).

\textbf{Output:} The substratum is a \(d = 3\) cubic lattice bijection
with: - \(K = 2d = 6\) internal components per site - Coupling matrix
eigenvalue multiplicities \((3, 2, 1)\) - Gauge group
\(\text{SU}(3) \times \text{SU}(2) \times \text{U}(1)\) - Three
generations of chiral fermions in the SM representations - One Higgs
doublet \((\mathbf{1}, \mathbf{2}, +1/2)\) - Anomaly-free hypercharges
\((Y_Q, Y_u, Y_d, Y_L, Y_e) = (1/6, 2/3, -1/3, -1/2, -1)\) -
\(\bar\theta = 0\) exactly

\textbf{Derivation.}

\begin{enumerate}
\def\labelenumi{(\alph{enumi})}
\item
  \emph{Dimension.} The coupling graph inherited from Stage 1 has
  polynomial growth exponent \(d\). Three independent forward filters
  from the empirical inputs converge on \(d = 3\): propagating gravity
  (E5) requires \(d \geq 3\) (the Weyl tensor vanishes for
  \(d \leq 2\)); stable matter (E6) requires \(d \leq 3\) (Coulomb
  instability of atoms for \(d \geq 4\), plus Tangherlini's
  gravitational orbit instability); boundary-entropy concordance (E7)
  gives \(\rho_s / \rho_{\text{crit}} = 2/(d-1) = 1\) at \(d = 3\)
  uniquely. The intersection of the first two filters already pins
  \(d = 3\); the third is a sharpening consistency check that the
  framework passes exactly. The integer-polynomial growth assumption
  (rather than fractal or exponential growth) is itself derived from
  self-consistency: exponential growth fails Lorentz-invariant
  dispersion (Jacobson's GR derivation requires it); fractal growth
  fails the cubic-lattice symmetry needed for the cubic-group
  decomposition in step (d) below. By Gromov's theorem and statistical
  isotropy (E4), the surviving graph class is quasi-isometric to
  \(\mathbb{Z}^d\) for integer \(d\); the three filters then select
  \(d = 3\). ({[}SM §3.2{]} presents this argument in full; the inputs
  E5, E6, E7 surfaced here make explicit the empirical content the
  filters require.)
\item
  \emph{Wave equation.} Center independence (A4), isotropy (E4), and
  linearity (A5) uniquely select the wave equation up to a
  multiplicative coupling constant \(\alpha\) ({[}SM §4.1, Theorem{]}):
  the unique second-order linear dynamics on a lattice that is
  translation-invariant, isotropic, and reversible has the form
  \(f = \alpha(x_1 + x_2 + \cdots + x_{2d}) \bmod q\) with propagation
  speed \(v = \alpha\). The constant \(\alpha = 1\) is fixed by
  relativistic causality with maximum signal speed \(c\) realized at the
  lattice cutoff.
\item
  \emph{Internal components.} Coupling-degree minimization applied to
  the wave equation gives \(K = 2d = 6\) uniquely (Theorem 6). This
  locks the absolute generation count at three (three spin-1/2 staggered
  tastes emerge from the \(K = 6\) minimum).
\item
  \emph{Gauge group.} Cubic-group decomposition of the \(K = 6\)
  components gives multiplicities \((3, 2, 1)\) (Theorem 7). The
  pointwise commutant of the cubic-group action on
  \(V_K = V_3 \oplus V_2 \oplus V_1\) is naturally
  \(U(3) \times U(2) \times U(1)\) (the full centralizer in \(GL(V_K)\)
  restricted to unitary transformations). The reduction to
  \(SU(3) \times SU(2) \times U(1)\) occurs because the overall \(U(1)\)
  phase of each block is gauge under generator (ii) of
  \(\mathcal{G}_{\text{sub}}\) (alphabet-size freedom \(q\) in §4): a
  global rescaling of the \(V_3\) block by a phase is indistinguishable
  from a global rescaling of \(\phi\), which is gauge by
  alphabet-freedom. The same applies to the \(V_2\) block. The
  singlet-block \(U(1)\) survives this stripping (it cannot be absorbed
  into alphabet-freedom because it acts non-trivially on the other
  blocks via the determinant constraint) and becomes the hypercharge
  \(U(1)\). Background independence (A6) then promotes the remaining
  global \(SU(3) \times SU(2) \times U(1)\) to local gauge invariance
  ({[}SM §3.1{]}).
\item
  \emph{Matter content.} Staggered reduction gives three degenerate
  spin-1/2 sectors (Theorems 8--10). Fermion embedding is completed via
  the link-carrier construction ({[}SM §4.7.1.2{]}). Partition-spinor
  identification (Theorem 12) and trace-out (Theorem 13) give chirality.
  Anomaly cancellation (Theorems 14--15) uniquely fixes the
  hypercharges.
\item
  \emph{Discrete symmetries.} T-invariance of the bijection \(\varphi\)
  combined with detailed balance in the emergent Hamiltonian forces
  \(\bar\theta = 0\) at all energy scales (Theorems 17--21).
\end{enumerate}

\textbf{Uniqueness at Stage 2.} Each step (a)--(f) is an
\emph{if-and-only-if}: the stated structure is the unique consistent
choice given the inputs. The full chain is therefore a composition of
unique selections, so the output is unique up to
\(\mathcal{G}_{\text{sub}}\).

\textbf{Status:} Theorem at the lattice level.

\hypertarget{stage-3-dynamics-emergent-fundamental-constants}{%
\subsubsection{3.4 Stage 3: Dynamics → emergent fundamental
constants}\label{stage-3-dynamics-emergent-fundamental-constants}}

\textbf{Inputs:} Stage 2 output.

\textbf{Output:} The emergent value of \(\hbar\) is
\(\hbar = c^3 \epsilon^2 / (4G)\) with \(\epsilon = 2 l_p\), where
\(l_p\) is the Planck length.

\textbf{Derivation.} The gap equation from thermal self-consistency of
the boundary layer ({[}GR §3{]}) gives the stated relation. The
derivation uses: existence of thermal equilibrium at the boundary layer
(a consequence of the C1--C3 dynamics from Stage 1), the boundary-only
dependence lemma ({[}GR §3.2{]}), detailed balance (consequence of
\(\varphi\) being a bijection). No new structural assumptions beyond
Stages 1--2.

\textbf{Uniqueness at Stage 3.} The gap equation admits a unique
solution under the stated inputs ({[}GR §3, Theorem{]}).

\textbf{Status:} Theorem.

\hypertarget{the-composite-theorem}{%
\subsubsection{3.5 The composite theorem}\label{the-composite-theorem}}

\textbf{Lemma 23.0} (Uniqueness preservation). \emph{Let
\((S_1, \varphi_1)\) and \((S_2, \varphi_2)\) both be reconstructions of
the same observed physics (E1--E7) satisfying the structural assumptions
(A1--A6). Then \((S_1, \varphi_1)\) and \((S_2, \varphi_2)\) are in the
same equivalence class \([(S, \varphi)]/\mathcal{G}_{\text{sub}}\).}

\emph{Proof.} By Stage 1, each \((S_i, \varphi_i, V_i)\) is a C1--C3
embedding of the observed QM. By the Stinespring uniqueness and the
\(\mathcal{G}_{\text{sub}}\) identification of Stage 1, the two
embeddings are related by generators (i) and (iii) of
\(\mathcal{G}_{\text{sub}}\) (state relabeling and deep-sector
enlargement).

By Stage 2, each triple is placed on a \(d = 3\) cubic lattice with
\(K = 6\) and the specific gauge structure. Any two lattices satisfying
A3--A6 and having the observed isotropy (E4) are related by generator
(iv) of \(\mathcal{G}_{\text{sub}}\) (graph isomorphism up to
statistical isotropy). The alphabet size freedom is generator (ii).

By Stage 3, each triple produces \(\hbar = c^3 \epsilon^2 / (4G)\). The
constant is the same in both reconstructions.

Therefore \((S_1, \varphi_1)\) and \((S_2, \varphi_2)\) are related by
the composition of generators (i)--(iv) of \(\mathcal{G}_{\text{sub}}\),
and hence are in the same equivalence class. \(\square\)

\textbf{Theorem 23} (Layered reconstruction). \emph{Let E1--E7
(empirical inputs: QM with Bell violations, finite boundary entropy,
spatial isotropy, propagating gravity, stable matter, boundary-entropy
concordance) and A1--A6 (structural assumptions: finiteness,
determinism, bounded coupling degree, center independence, linearity,
background independence) hold. Then the equivalence class
\([(S, \varphi)]/\mathcal{G}_{\text{sub}}\) is uniquely determined,
conditional on ETH in the hidden sector for the C2 necessity direction
(Stage 1): a finite set with a bijection of bounded coupling degree and
statistical isotropy, with \(d = 3\), \(K = 6\), coupling matrix
eigenvalue multiplicities \((3, 2, 1)\), gauge group
\(\text{SU}(3) \times \text{SU}(2) \times \text{U}(1)\), three chiral
generations, one Higgs doublet, anomaly-free hypercharges,
\(\bar\theta = 0\), and \(\hbar = c^3 \epsilon^2/(4G)\).}

\emph{Proof.}

\begin{enumerate}
\def\labelenumi{(\arabic{enumi})}
\item
  Stage 1 (§3.2) establishes existence and uniqueness (up to
  \(\mathcal{G}_{\text{sub}}\)) of a C1--C3 embedding from E1--E3 +
  A1--A2.
\item
  Stage 2 (§3.3) takes the Stage 1 output plus E4 + A3--A6 and derives
  the specific lattice, dimension, gauge group, matter content, and
  discrete symmetries, with uniqueness at each step.
\item
  Stage 3 (§3.4) takes the Stage 2 output and derives the emergent
  constant \(\hbar\), with uniqueness.
\item
  Lemma 23.0 establishes that any two reconstructions consistent with
  the inputs are in the same equivalence class, completing the
  uniqueness claim.
\end{enumerate}

The existence direction is the composition of Stages 1--3. The
uniqueness direction is Lemma 23.0. Together they establish the
bidirectional correspondence. \(\square\)

\hypertarget{remarks}{%
\subsubsection{3.6 Remarks}\label{remarks}}

\textbf{Remark} (Evidential weight of the reconstruction). The
reconstruction uniquely derives
\(\text{SU}(3) \times \text{SU}(2) \times \text{U}(1)\) with
multiplicities \((3, 2, 1)\), three chiral generations, one Higgs
doublet \((1, 2, +1/2)\), unique hypercharges, and \(\bar\theta = 0\).
Because the derivation is unique --- no free parameters, no
model-building choices, no landscape of alternatives --- every
independent experimental confirmation of the Standard Model's gauge
structure is retroactive evidence for the derivation chain that produces
it. The Standard Model is the most precisely tested mathematical
structure in the history of science (\(g - 2\) of the electron confirmed
to \(10^{-12}\), electroweak precision observables to \(10^{-5}\), QCD
cross-sections across decades of energy). This experimental record
transfers in full to the framework, because uniqueness of the derivation
makes the relationship between evidence and theory transitive: data
confirms structure, structure is uniquely derived, therefore data
confirms the derivation.

\textbf{Remark} (Scope of uniqueness). The theorem establishes
uniqueness within the class of candidate substrates satisfying A1--A6.
It does \textbf{not} exclude:

\begin{itemize}
\tightlist
\item
  Intrinsically continuum theories (violate A1)
\item
  Intrinsically stochastic theories (violate A2)
\item
  Theories with unbounded coupling degree (violate A3)
\item
  Theories violating center independence, linearity, or background
  independence (violate A4--A6)
\end{itemize}

The framework's claim is that within the class of finite deterministic
bounded-coupling systems with the structural assumptions, the substratum
is unique. A separate argument would be required to rule out
alternatives outside this class. The framework's defense of A1--A6 (in
{[}SM §2{]} and elsewhere) is that they are natural given the empirical
inputs E1--E7, but this defense is itself an argument, not a theorem.

\textbf{Remark} (Hypothesis dependencies). Each structural assumption
can in principle be weakened or replaced:

\begin{itemize}
\tightlist
\item
  A1 (finiteness) is supported by E3; alternative would be an
  effective-finiteness argument with continuum UV completion, which the
  framework does not develop.
\item
  A2 (determinism) is the key ontological commitment; stochastic
  alternatives would produce different reconstructions.
\item
  A3 (bounded coupling) is conventional for physical systems; long-range
  couplings would give different dimensional structure.
\item
  A4 (center independence) rules out preferred-frame theories.
\item
  A5 (linearity) is the strongest restriction; nonlinear wave equations
  on finite lattices are a separate class.
\item
  A6 (background independence) is standard for gauge theories.
\end{itemize}

Weakening any \(A_i\) does not merely enlarge the equivalence class; it
potentially produces \emph{different} reconstructions that may not agree
with observation. The framework's claim is that E1--E7 + A1--A6 is a
consistent set of inputs producing our observed physics; alternative
input sets are not investigated.

\textbf{Remark} (Continuum extension). The lattice-level predictions are
the framework's primary claims --- the lattice is the fundamental
description, not an approximation to a continuum theory. The structural
results (gauge group, representations, generation count,
\(\bar\theta = 0\)) are algebraic/topological properties of the lattice
theory and are exact. Quantitative observables (scattering amplitudes,
mass ratios) are compared to experiment through continuum perturbation
theory; the lattice-continuum discrepancy at any experimentally
accessible energy \(E\) is suppressed by
\((E\epsilon/\hbar c)^2 = (E/M_{\text{Pl}})^2\), which is
\(\sim 10^{-32}\) at LHC energies. The rigorous proof that lattice
Yang-Mills defines a continuum limit with a mass gap (a Clay Millennium
Prize problem) is an open problem in mathematics, but it is not required
by the OI framework: the lattice theory is complete, and its predictions
are lattice-exact.

\textbf{Remark} (Finite observation sets vs.~idealized empirical
inputs). The empirical inputs E1--E7 are stated as structural facts
about observed physics --- the full apparatus of QM, Bell violations
attaining Tsirelson's bound, the holographic entropy bound, exact
rotational invariance, propagating gravitational waves, stable matter on
all observed scales, and spatial-curvature concordance. Operationally,
no finite observation set fully establishes any of these; finite
measurements support them only in the sense of being consistent with the
data so far. The reconstruction theorem holds for the empirical inputs
\emph{as stated} --- the idealized infinite-data regime --- and proves
uniqueness of the equivalence class under that regime. Reconstruction
from finite observation sets is a related but distinct question: with
finite data, multiple equivalence classes may all be consistent, with
the data discriminating among them only weakly. The framework's
operational answer to the finite-data question is the cumulative weight
of the quantitative predictions of {[}SM §7{]} --- twenty-two structural
retrodictions of gauge couplings, CKM/PMNS angles, mass ratios, and the
Koide relation, each at \(\lesssim 1\sigma\) --- which collectively rule
out reconstructions that would have produced different values. Theorem
23 establishes idealized uniqueness; the {[}SM §7{]} prediction set
establishes the operational uniqueness that a skeptical reader can check
against data.

The reconstruction establishes a bidirectional correspondence:

\[\text{Observed physics (E1–E7)} \; + \; \text{Structural assumptions (A1–A6)} \quad \longleftrightarrow \quad [(S, \varphi)] / \mathcal{G}_{\text{sub}}\]

The mathematical structure and the physics determine each other up to
gauge equivalence, given the structural assumptions. The distinction
between ``mathematics describes reality'' and ``mathematics is reality''
has no empirical content --- it is itself gauge, provably undecidable by
measurement (cf.~Theorem 24). This reframes Wigner's puzzle: the
``unreasonable effectiveness'' of mathematics is a theorem modulo the
structural assumptions, not a mystery.

\textbf{Remark} (ER=EPR at substratum and emergent levels). The
Maldacena-Susskind ER=EPR conjecture {[}3{]} proposes that any two
entangled systems are connected by an Einstein-Rosen bridge, with the
strong form asserting that entanglement \emph{is} a non-traversable
wormhole --- the same ontological structure under two descriptions. The
framework here engages this conjecture at two distinct levels.

At the \textbf{emergent level}, the conjecture's \emph{weak} form is
reproduced: any pair of entangled subsystems in the emergent QM
description shares boundary modes via the partition trace-out, and the
holographic dictionary maps shared boundary modes to a connecting
geometric structure in the emergent gravitational description. This
recovers the ER=EPR weak form --- entangled systems are connected by a
wormhole-like geometric object in the emergent theory --- as a
structural consequence of the framework's holographic architecture
rather than as an independent postulate.

At the \textbf{substratum level}, the conjecture's \emph{strong} form
(entanglement is identical to wormhole geometry, not merely correlated
with it) is predicted \emph{false} at the representative level. Two
distinct partition representatives can produce identical emergent EPR
correlations while having different substratum representative structures
--- there is no faithful bijection from EPR pairs to substratum
geometric features. The strong identity holds at most up to gauge
equivalence on the substratum side. Whether a gauge-orbit reformulation
of strong ER=EPR --- ``EPR-equivalence-classes correspond to
ER-equivalence-classes'' --- is true in the framework remains open; the
test would be whether the substratum's gauge group
\(\mathcal{G}_{\text{sub}}\) acts transitively on the set of substratum
representatives that produce a given EPR correlation pattern.

This produces a sharp empirical signature: the framework's operational
predictions agree with weak ER=EPR (which would also be predicted by
AdS/CFT, by quantum error correction in holography, and by other
emergent-geometry approaches) but disagree with strong ER=EPR at the
representative level. Distinguishing these experimentally would require
probing substratum-scale features, which is currently beyond reach. The
disagreement is therefore a structural prediction rather than a testable
one at present, but it concretely positions the framework relative to
the broader ER=EPR program.

\begin{center}\rule{0.5\linewidth}{0.5pt}\end{center}

\hypertarget{the-substratum-gauge-group}{%
\subsection{4. The Substratum Gauge
Group}\label{the-substratum-gauge-group}}

The equivalence relation \(\sim\) in the reconstruction theorem has a
precise structure. Define: \((S, \varphi) \sim (S', \varphi')\) if the
two systems produce identical emergent physics --- the same transition
probabilities \(T_{ij}(t)\), the same emergent Hamiltonian (up to
D-gauge), and the same \(\hbar\) --- for all partitions of the same
structural class. The set of transformations mapping \((S, \varphi)\) to
an equivalent \((S', \varphi')\) is a group --- the \emph{substratum
gauge group} \(\mathcal{G}_{\text{sub}}\).

\textbf{Theorem 24} (Generators of the substratum gauge group).
\emph{\(\mathcal{G}_{\text{sub}}\) contains at least four independent
families of transformations:}

\emph{(i) State relabeling.} For any bijection \(\sigma: S \to S\), the
conjugate
\((S, \sigma \circ \varphi \circ \sigma^{-1}) \sim (S, \varphi)\). The
transition probabilities depend on the coupling structure of
\(\varphi\), not on which labels are attached to which states. This is
an \(|S|!\)-element subgroup --- vastly larger than any gauge group in
the Standard Model.

\emph{(ii) Alphabet change.} Replacing the local state space
\(\mathbb{Z}/q\mathbb{Z}\) with \(\mathbb{Z}/q'\mathbb{Z}\) for any
\(q' \geq 2\), while preserving the coupling graph and dynamics class,
leaves all observables unchanged ({[}SM, §2.7{]}). This family is
parametrized by all integers \(q \geq 2\).

\emph{(iii) Deep-sector enlargement.} Adjoining additional degrees of
freedom to \(\mathcal{C}_D\) (the deep hidden sector beyond the boundary
layer), with arbitrary dynamics satisfying \(\tau_B^D \gg \tau_S\), does
not change the emergent description. The boundary-only dependence lemma
{[}GR, §3.2{]} proves
\(T_{ij}(t) = T_{ij}^{(B)}(t) + \mathcal{O}(t/\tau_B)\): observables
depend only on \(\mathcal{C}_V \times \mathcal{C}_B\). The deep sector
may be finite of any size, or infinite.

\emph{(iv) Graph isomorphism (up to statistical isotropy).} Two coupling
graphs \(G_\varphi\) and \(G_{\varphi'}\) that are quasi-isometric with
the same polynomial growth exponent \(d\), the same spectral properties,
and the same statistical isotropy at large scales produce the same
emergent physics. The regular cubic lattice \(\mathbb{Z}^3\) and any
bounded-degree random graph with \(d = 3\) polynomial growth and
statistical isotropy are gauge-equivalent.

\emph{Proof.} Each generator preserves all inputs to the derivation
chain. (i): conjugation preserves the coupling graph \(G_\varphi\) up to
relabeling, hence all graph-dependent quantities (area law, dispersion,
dimension, eigenvalue multiplicities). (ii): {[}SM, §2.7{]} proves
\(q\)-independence of every prediction. (iii): the boundary-only
dependence lemma gives the result directly. (iv): the derivation chain
uses only statistical properties of \(G_\varphi\) (dimension via
Myrheim-Meyer, isotropy, bounded degree), not the specific graph.
\(\square\)

\textbf{Completeness of the generators.} The four families exhaust
\(\mathcal{G}_{\text{sub}}\). The proof proceeds by showing that any
observables-preserving transformation decomposes into (i)--(iv).

Let \(g: (S, \varphi) \to (S', \varphi')\) preserve all observables
(transition probabilities \(T_{ij}(t)\), emergent Hamiltonian up to
D-gauge, and \(\hbar\)). By assumption, \(g\) preserves the
visible-sector CPTP channel \(\Phi\). Stinespring's uniqueness theorem
{[}Main, §3.2{]} constrains the freedom: any two dilations of the same
channel on the same visible Hilbert space differ by a partial isometry
on the hidden sector. At the substratum level, this partial isometry is
a relabeling of hidden-sector states (generator (i) restricted to
\(\mathcal{C}_H\)) when the hidden sectors have equal cardinality, or a
composition of relabeling and deep-sector enlargement/reduction
(generators (i) + (iii)) when they differ in size. The alphabet may
change freely (generator (ii), by {[}SM, §2.7{]}). The coupling graph
\(G_{\varphi'}\) must have the same statistical properties as
\(G_\varphi\) --- dimension, spectral dimension, isotropy --- since
these determine the derivation chain's outputs; two such graphs are
related by generator (iv). These four degrees of freedom ---
hidden-sector relabeling, deep-sector size, alphabet, and graph
statistical class --- exhaust the free parameters in the reconstruction,
so the four generators span \(\mathcal{G}_{\text{sub}}\). \(\square\)

Three candidate fifth families are explicitly subsumed:

\begin{itemize}
\tightlist
\item
  \emph{Time reversal} (\(\varphi \to \varphi^{-1}\)): the wave
  equation's T-invariance gives
  \(\varphi^{-1} = T \circ \varphi \circ T^{-1}\) where \(T\) is the
  phase-space layer swap \((x(t), x(t+1)) \mapsto (x(t+1), x(t))\). This
  is generator (i) with \(\sigma = T\).
\item
  \emph{Hidden-sector dynamical reparametrization} (beyond enlargement):
  by Stinespring uniqueness, any two same-channel dilations of equal
  hidden-sector dimension differ by a hidden-sector unitary ---
  generator (i) restricted to \(\mathcal{C}_H\).
\item
  \emph{Visible-sector emergent global phase} (\(U \to e^{i\theta}U\)):
  at the substratum level \(S\) is a finite set with no complex
  structure; the emergent phase is trivially the identity on
  \((S, \varphi)\).
\end{itemize}

\textbf{The gauge hierarchy.} Three levels of gauge symmetry appear in
the framework, each projecting onto the next through the trace-out:

\emph{Level 3 (substratum):} \(\mathcal{G}_{\text{sub}}\) acts on
\((S, \varphi)\) before the trace-out. It is the largest gauge group and
includes transformations with no analog in the emergent description
(deep-sector enlargement, alphabet change).

\emph{Level 2 (emergent QFT):}
\(\text{SU}(3) \times \text{SU}(2) \times \text{U}(1)\) is the commutant
of the coupling matrix \(M\), acting on the emergent fields. It is the
image of \(\mathcal{G}_{\text{sub}}\) restricted to transformations that
permute internal components within the eigenspaces of \(M\).

\emph{Level 1 (emergent Hamiltonian):} The D-gauge \(H \to DHD^\dagger\)
with \(D\) a diagonal unitary, acting on the emergent Hamiltonian within
the emergent QM. It is the residual freedom after all
transition-probability data has been extracted.

Each level is contained in the one above: Level 1 \(\subset\) Level 2
\(\subset\) Level 3. The trace-out projects Level 3 onto Level 2 (the SM
gauge group is the shadow of \(\mathcal{G}_{\text{sub}}\) visible to the
emergent QFT), and restricting to the Hamiltonian projects Level 2 onto
Level 1.

\textbf{Remark.} The substratum gauge group is not a symmetry of a
Lagrangian or an action --- no Lagrangian exists at the substratum
level. It is a symmetry of the \emph{equivalence class of substrata},
defined by the condition that all observables are preserved. The
emergent gauge symmetries (Levels 1 and 2) are Lagrangian symmetries in
the standard sense, derived from the substratum through the trace-out.

\begin{center}\rule{0.5\linewidth}{0.5pt}\end{center}

\hypertarget{synthesis-three-projections-of-s-ux3c6}{%
\subsection{5. Synthesis: Three Projections of (S,
φ)}\label{synthesis-three-projections-of-s-ux3c6}}

The reconstruction theorem (§3) establishes that observed physics ---
quantum mechanics with Bell violations, finite boundary entropy, and
spatial isotropy --- uniquely determines the equivalence class
\([(S, \varphi)]/\mathcal{G}_{\text{sub}}\) at the lattice level. The
substratum gauge group (§4) makes the equivalence relation precise and
identifies the Standard Model gauge group as the visible-sector shadow
of \(\mathcal{G}_{\text{sub}}\). This section makes the synthesis claim
explicit: quantum mechanics, general relativity, and the arrow of time
are not three independent theories but three projections of the same
finite deterministic construction.

\hypertarget{the-three-projections}{%
\subsubsection{5.1 The three projections}\label{the-three-projections}}

The framework's derivations apply the same trace-out machinery to the
same \((S, \varphi)\) along three projections: a visible-sector
projection that produces emergent QM ({[}Main{]}) and the Standard Model
({[}SM{]}); a boundary-thermodynamic projection that produces general
relativity ({[}GR{]}); and an arrow-of-time projection (§5.4) that
produces the emergent cosmological, thermodynamic, memory, and
measurement arrows from a single structural clock and the C1--C3
observer-selection theorem of {[}Main, §4.6{]}.

\textbf{Visible-sector projection.} At the level of the embedded
observer's epistemic access, the trace-out over the hidden sector
produces unitary quantum mechanics with the wave function, Born rule,
and Schrödinger evolution as derived structures (rather than postulated
ones). The proof in {[}Main, §3{]} establishes this as a theorem: under
the three structural conditions C1 (non-trivial coupling), C2 (slow
bath), and C3 (high-capacity hidden sector), the embedded observer's
compressed description is necessarily quantum mechanical, with the
conditions both necessary and sufficient. Applied to a cubic lattice
with the wave equation as substratum dynamics --- itself uniquely
selected among second-order reversible nearest-neighbor dynamics by
center independence, isotropy, and linearity ({[}SM, §4.1{]}) --- this
projection produces the Standard Model gauge group
\(\mathrm{SU}(3) \times \mathrm{SU}(2) \times \mathrm{U}(1)\) (as the
commutant of the coupling matrix, {[}SM, §§4.4--4.6{]}), three chiral
generations and the Higgs as a composite scalar in the singlet staggered
taste ({[}SM, §4.7{]}, Theorems 8--11), anomaly-free hypercharges
({[}SM, §4.9{]}, Theorems 14--15), and \(\bar\theta = 0\) at all energy
scales ({[}SM, §5{]}, Theorems 17--21). Twenty-two quantitative
predictions follow ({[}SM, §7{]}), including the Cabibbo angle from a
single Brillouin-zone distance, the Koide angle from a cubic-group
quadratic Casimir, all three PMNS angles within \(1.1\sigma\) (vs NuFIT
6.0), six fermion masses from one empirical input to better than
\(1\%\), and all three SM gauge couplings at \(M_Z\) ({[}SM, §6.3{]}).

\textbf{Boundary-thermodynamic projection.} At the level of the
partition boundary, classical horizon thermodynamics --- the
temperature, entropy, and dynamical evolution of the causal horizon as a
gravitating object --- produces general relativity through Jacobson's
thermodynamic argument: the Clausius relation \(dE = T \, dS\) applied
at local causal horizons yields Einstein's equations with the same
structural inputs that give the visible-sector projection its form.
Applied to the cosmological horizon as a causal partition ({[}GR,
§2{]}), this projection determines the emergent action scale
\(\hbar = c^3 (2 l_p)^2 / (4G)\) from thermal self-consistency ({[}GR,
§3{]}), fixes the discreteness scale \(\epsilon = 2 l_p\) as the unique
simultaneous solution ({[}GR, §4{]}), derives the Bekenstein-Hawking
entropy \(S = A/(4 l_p^2)\) including the \(1/4\) coefficient by direct
mode counting ({[}GR, §5{]}) --- recently confirmed at \(99.999\%\)
confidence by GW250114 --- and dissolves the cosmological constant
problem by identifying the quantum vacuum energy and the effective
\(\Lambda\) as properties of logically distinct levels of description
rather than commensurable quantities whose \(10^{122}\) ratio requires
cancellation ({[}GR, §6{]}). The framework predicts dynamical dark
energy in Type II running-vacuum functional form as a structural
prediction (magnitude not derived at theorem level), consistent with the
Bertini et al.~(2025) direct DESI-DR2-era fit at \(2\sigma\) (which is
consistent with the ΛCDM limit); a \textasciitilde{}\(95\%\) dark sector
corollary as a structural feature rather than a new particle species;
and a MOND acceleration scale \(a_0 = cH/6\) from entropy displacement
at the boundary ({[}GR, §7{]}).

The three projections share the same source (\((S, \varphi)\)), the same
trace-out machinery (marginalization over the hidden sector under
conditions C1--C3), and the same structural inputs (the partition
geometry, the boundary entropy, the substratum gauge group). They differ
only in which structure of the embedded observer's description they
derive: the visible-sector projection (§5.1, §5.3) works on the bulk
dynamics and produces emergent QM and the Standard Model; the
boundary-thermodynamic projection (§5.1) works on the classical thermal
data at the partition boundary and produces general relativity; the
arrow-of-time projection (§5.4) works on the horizon-clock parameter
along observer worldlines and produces the cascade of emergent arrows.

\hypertarget{the-qm-gr-incompatibility-as-a-category-error}{%
\subsubsection{5.2 The QM-GR incompatibility as a category
error}\label{the-qm-gr-incompatibility-as-a-category-error}}

The conventional formulation of the QM-GR unification problem treats
quantum mechanics and general relativity as two competing theories in
the same logical category --- two attempts to describe the same regime
--- and asks how to merge them into a single mathematically consistent
framework. Every program of unification proceeds from this assumption:
quantum gravity programs treat the metric as a quantum field to be
quantized; geometric programs treat the wave function as a structure on
the spacetime manifold; emergent programs treat one of the two as
derived from a deeper substrate that recovers the other. The persistent
failure of all three approaches to produce a quantitatively confirmed
unification suggests that the underlying assumption is wrong.

In the present framework, the assumption is wrong because quantum
mechanics and general relativity occupy different positions in the
trace-out hierarchy. They are not two theories of the same regime at all
--- they are two projections of the same construction, viewed at
different levels. The visible-sector projection that produces QM and the
boundary-thermodynamic projection that produces GR are not in tension
because they refer to \emph{different objects in the construction}. The
wave function is a property of the embedded observer's compressed
description of the visible sector; the metric is a property of the
classical dynamics of the partition boundary. These are not two
descriptions of the same physical system at different levels of
approximation. They are descriptions of two different substructures of
the same total object \((S, \varphi)\).

The analogy that captures the structure most clearly is the one between
a thermodynamic and a statistical-mechanical description of a gas. The
thermodynamic description deals with pressure, temperature, and entropy;
the statistical-mechanical description deals with particle positions,
momenta, and microstate counting. These are not two competing theories
of the gas, and the question ``how do we unify thermodynamics and
statistical mechanics?'' is not a coherent question --- they describe
different things about the same system, and the relationship between
them is one of projection (statistical mechanics produces thermodynamic
quantities by averaging) rather than unification. The QM-GR relationship
in the present framework is structurally similar: GR describes the
thermodynamic (boundary, classical, deterministic) projection of
\((S, \varphi)\), and QM describes the statistical (visible-sector,
compressed, stochastic-emergent) projection. The two are related by a
definite procedure (the trace-out under conditions C1--C3), and the
question ``how do we unify them?'' is dissolved rather than answered.

This is not the same as saying that quantum mechanics is ``more
fundamental'' than general relativity, or vice versa. Both are emergent.
Both depend on the trace-out and the partition. Both are derived rather
than fundamental. The fundamental object is \((S, \varphi)\) --- a
finite set with a deterministic bijection --- and neither QM nor GR is a
feature of \((S, \varphi)\) itself. They are features of how an embedded
observer or a horizon-bounded thermodynamic regime sees \((S, \varphi)\)
from inside.

\hypertarget{the-three-level-hierarchy-made-explicit}{%
\subsubsection{5.3 The three-level hierarchy made
explicit}\label{the-three-level-hierarchy-made-explicit}}

The substratum gauge group (§4) provides the structural framework for
the synthesis claim at the level of the two gauge-bearing projections
(the arrow-of-time projection is developed in §5.4 and has no
gauge-hierarchy analog). At the substratum level, the only object is the
bijection and its gauge group \(\mathcal{G}_{\text{sub}}\). The
trace-out projects this onto two derived structures simultaneously: the
emergent quantum field theory at the visible-sector level, with the
Standard Model gauge group as the commutant of the coupling matrix
({[}SM, §4.4{]}, Theorem 5) and the Standard Model representations as
the cubic decomposition of the link directions ({[}SM, §4.6{]}, Theorem
7); and the classical horizon thermodynamics at the boundary level, with
the metric, the surface gravity, the entropy, and the temperature
determined by the partition geometry and the area-law theorem ({[}GR,
§3{]}). Both derived structures inherit the framework's structural
inputs from the substratum, and both are exact rather than approximate
at the lattice level.

The three-level gauge hierarchy makes this exact. At Level 3 (the
substratum), \(\mathcal{G}_{\text{sub}}\) acts on \((S, \varphi)\)
before any trace-out and includes transformations with no analog in the
emergent description (state relabeling, alphabet change, deep-sector
enlargement, graph isomorphism up to statistical isotropy). At Level 2
(the emergent QFT), the Standard Model gauge group acts on the emergent
fields as the commutant of the coupling matrix, and is the image of
\(\mathcal{G}_{\text{sub}}\) restricted to transformations that permute
internal components within the eigenspaces of \(M\). At Level 1 (the
emergent Hamiltonian), the diagonal-unitary D-gauge acts on the emergent
Hamiltonian within the emergent QM and is the residual freedom after all
transition-probability data has been extracted. Each level is contained
in the one above, and the trace-out projects each onto the next.

The boundary-thermodynamic projection that produces GR is a parallel
structure at Level 2 --- but acting on the \emph{partition} rather than
on the internal field content. Where the emergent QFT is the trace-out's
effect on the bulk dynamics, the classical horizon thermodynamics is the
trace-out's effect on the partition geometry itself. Both are at Level
2. Both descend from Level 3. The framework's claim is that this
structural relationship is exact: the SM gauge group and Einstein's
equations are not two independent theoretical inputs but two co-derived
structures, both consequences of the same underlying \((S, \varphi)\)
and the same trace-out under the same conditions C1--C3.

\hypertarget{the-arrow-of-time-as-a-third-projection}{%
\subsubsection{5.4 The arrow of time as a third
projection}\label{the-arrow-of-time-as-a-third-projection}}

After Main §4.6 (structural observer-selection theorem), the arrow of
time joins quantum mechanics and general relativity as a third emergent
structure that the framework derives from the same underlying
\((S, \varphi)\) and the same trace-out machinery. Before §4.6, the
arrow of time was treated as external to the derivation chain --- a past
hypothesis or an anthropic selection, imposed on the framework rather
than produced by it. With §4.6 in place, the observer's confinement to
the non-equilibrium phase of \((S, \varphi)\) is a structural fact, and
every arrow the observer sees descends from this confinement together
with the horizon growth that the accelerating phase produces. The third
projection has the same form as the first two: a feature of how the
embedded observer sees \((S, \varphi)\) from inside, rather than a
feature of \((S, \varphi)\) itself.

\textbf{Horizon growth as the structural clock.} The framework's primary
arrow is the cosmological horizon entropy
\(S_{\rm dS}(t) = A(t)/(4 l_p^2)\), monotonically increasing along any
observer worldline in the \(\Lambda\)-dominated phase. The monotonicity
is a feature of the causal partition rather than of any local frame:
every observer in the accelerating phase agrees that \(S_{\rm dS}\)
increases, just as every observer in a Schwarzschild geometry agrees on
the exterior direction. The horizon clock is observer-independent,
continuous, and monotonic for as long as \(\ddot a > 0\), which the
framework identifies as the regime of interest for the derivations of
{[}GR{]}. At the substrate level, \((S, \varphi)\) has no preferred
direction: \(\varphi\) and \(\varphi^{-1}\) are equally valid as
forward-time maps. At the observer level, the horizon clock picks out a
direction, and all other arrows align with it.

\textbf{The layered cascade.} Four derived arrows align with the horizon
clock through a structural cascade. \emph{Cosmological:}
\(S_{\rm dS}(t)\) monotonic in the \(\Lambda\)-dominated phase, as
above. \emph{Thermodynamic:} via Jacobson's coupling
\(\rho_s = \rho_{\rm crit}\) ({[}GR, §3{]}), the entropy of any
horizon-bounded system increases with the observer's clock --- the
thermodynamic arrow descends from the structural arrow without an
independent postulate. \emph{Memory:} C2 (the slow bath) implies the
observer's hidden sector preserves correlations forward along the clock
rather than backward, because the spectral gap of the hidden-sector
dynamics is uniform in the two directions but the observer's measurement
protocol commits to one. \emph{Measurement:} the observer records
``past'' outcomes and prepares ``future'' states --- an asymmetry
inherited from the memory layer, not added on top of it. The chain
grounds in one structural fact (horizon growth) and one structural
constraint (C1--C3 observer selection via Main §4.6). No additional
axiom is needed.

\textbf{What the framework derives, and the directional-language
subtlety.} The framework's quantitative content at the level of Main
§4.6 is \emph{symmetric} in time direction --- the spectral-gap bound
\(\tau_B(V) \leq \tau_{\rm mix}(\Sigma_{\rm loc}(V)) + O(\epsilon)\)
applies equally in both directions along any observer worldline. What is
\emph{imposed} by the observer, not derived, is the convention of which
direction along the horizon-clock parameter to call ``future.'' The two
directions are gauge-equivalent at the substrate level
(\(\varphi \leftrightarrow \varphi^{-1}\)), and the observer's choice of
convention is the residual freedom after all physical content has been
extracted. The framework therefore derives the existence of a monotonic
parameter along observer worldlines, not a substrate-level arrow of
time. This distinction answers Ellis's objection (a bijective substratum
cannot carry an arrow of time) by agreeing with the objection's premise
while deriving what the observer actually sees from the structure the
substratum does have.

\textbf{Scope note: causal topology as load-bearing.} The memory-arrow
step uses ``the observer's worldline has a causal topology'' as the
carrier of the asymmetry. This is not a new axiom --- {[}Main, §1.3{]}
already assumes a causal topology on the partition --- but it is
load-bearing for the cascade and is noted explicitly here rather than
treated as obvious. A framework in which the observer's worldline had no
causal topology could not run the memory-arrow step, and the cascade
would halt at the thermodynamic layer.

\textbf{CP violation as the residual Ellis-class objection.} The layered
cascade closes the Ellis-axis objection for the arrow of time but does
not automatically close the related question of CP violation. The
substratum is T-symmetric (load-bearing for \(\bar\theta = 0\), {[}SM,
§5{]}), and observed CP violation in CKM and kaon/B systems is a
flavor-basis feature of the specific bijection rather than a consequence
of the cascade. The framework treats flavor-sector CP phases as
solution-specific inputs --- see {[}SM, §5.4{]} for the scope statement
and §6.4 ``what does not falsify'' for the commitment that the framework
makes no prediction of \(\delta_{\text{CKM}}\),
\(\delta_{\text{PMNS}}\), or the Wolfenstein \((\rho, \eta)\).

\hypertarget{what-the-synthesis-does-and-does-not-claim}{%
\subsubsection{5.5 What the synthesis does and does not
claim}\label{what-the-synthesis-does-and-does-not-claim}}

The synthesis claim is structural and architectural rather than
empirical. It does not claim that the framework predicts gravitational
phenomena that conventional QFT plus GR cannot reproduce --- most of
{[}GR{]}'s predictions (the BH area law, the basic properties of horizon
thermodynamics, the broad shape of dark energy phenomenology) are also
recovered by other approaches. The framework's empirical content is
concentrated in {[}SM{]} (the SM derivation, twenty-two quantitative
predictions) and in the specific GR-side predictions of {[}GR{]} (Type
II RVM dark energy as a structural-form prediction, MOND \(a_0 = cH/6\),
dark sector concordance), which collectively distinguish the framework
from standard \(\Lambda\)CDM plus the Standard Model.

What the synthesis claim \emph{does} establish is that those two papers
are not independent. The same construction that produces the SM in
{[}SM{]} produces the gravitational sector in {[}GR{]}, and the same
trace-out that gives the SM gauge group as the commutant of the coupling
matrix ({[}SM, §4.4{]}, Theorem 5) gives the BH entropy as the boundary
mode count ({[}GR, §5{]}). This is not a claim about new gravitational
phenomena. It is a claim about the structural relationship between two
derivations that, in the conventional picture, are independent and
incommensurable --- and that, in the present framework, are derived from
the same object by the same procedure under the same conditions.

This matters for two reasons. First, it removes the QM-GR unification
problem from the active research agenda by dissolving rather than
solving it: the problem was based on a category error, and the category
error is fixed by the substratum-level construction developed here.
Second, it provides an explanation for why the SM has the gauge group it
does --- namely, that the SM gauge group is the visible-sector shadow of
the substratum gauge group, with no choice of model and no landscape of
alternatives to select among. The SM is not one possibility among many
but the unique consequence of the framework's structural inputs.

\begin{center}\rule{0.5\linewidth}{0.5pt}\end{center}

\hypertarget{discussion}{%
\subsection{6. Discussion}\label{discussion}}

The reconstruction theorem identifies the question ``is mathematics
describing reality, or is it reality?'' as gauge in the precise sense
established by §4 --- provably undecidable by any measurement --- and
reframes Wigner's puzzle of the unreasonable effectiveness of
mathematics as a theorem rather than a mystery. The ontological
hierarchy makes explicit that space, time, matter, and energy are
derived rather than fundamental concepts. And the measurement problem
dissolves once the wave function is recognized as a derived object
rather than a component of the underlying reality. The pre-registration
of falsification conditions, finally, is the empirical counterpart of
the methodological discipline developed in §§3--5: a framework that
derives rather than posits must be willing to specify what would
invalidate the derivation.

\hypertarget{ontic-structural-realism-on-the-equivalence-class}{%
\subsubsection{6.1 Ontic structural realism on the equivalence
class}\label{ontic-structural-realism-on-the-equivalence-class}}

The metaphysical commitment introduced in §1.1 can now be stated with
the full technical machinery in place. The framework's claim is ontic
structural realism applied to
\([(S, \varphi)] / \mathcal{G}_{\text{sub}}\): what exists at the
substratum-class level is the equivalence class of bijections, uniquely
determined by observations plus structural assumptions (Theorem 23),
with individual substratum representatives related by the substratum
gauge group (Theorem 24). The class is real; specific substrates are
representations.

This is the strongest claim the framework's theorems license at this
level without overclaiming. Two alternative framings merit explicit
comparison.

\emph{Specific-substratum realism} --- the stronger commitment that a
particular bijection \(\varphi\) is the one our universe is in --- is
forbidden by Theorem 24. The four generators of
\(\mathcal{G}_{\text{sub}}\) relate representatives that produce
identical observables, so no measurement can discriminate among them;
specific-substratum realism commits to facts the framework's own gauge
structure classifies as unobservable.

\emph{Pure explanatory economy} --- the weaker position that the
framework is simply more economical than standard QM (fewer postulates,
more derivations) without ontic commitment --- undersells what Theorem
23 establishes. If observations plus structural assumptions uniquely
determine the equivalence class, the class is not one possibility among
many but the unique structure consistent with the evidence. The
explanatory economy follows from this uniqueness rather than
substituting for it.

The wave function, Hilbert space, emergent gauge group, and metric are
features of representations --- of the emergent description of a
specific partition of a specific substratum --- not of the equivalence
class itself. At the structural level there is only the bijection and
its symmetries; no quantum, no classical, no metric, no wave function.
This is the structural reason the QM-GR incompatibility dissolves: both
descriptions are emergent from the same structural object at different
levels of projection (§5). Asking whether the wave function is ``real''
is then asking about a feature of the representation, and the question
has a two-level answer --- the wave function is real at the emergent
level, gauge at the substratum level.

The question ``is mathematics describing reality, or is it reality?''
is, on this reading, provably gauge in the same technical sense that the
alphabet size \(q\) and the relabeling of hidden-sector states are
gauge. What the framework establishes is that the structural content is
determinate and unique; the metaphysical attitude one takes toward that
content is itself structural information the framework does not carry,
and therefore information no measurement can determine. Wigner's
``unreasonable effectiveness of mathematics'' becomes a theorem modulo
the structural assumptions rather than a mystery: physics is structural,
mathematics is the language of structure, and the effectiveness is the
reconstruction theorem.

\textbf{Multi-level structural realism.} The structural-realism
commitment is not single-level. {[}Structure{]} establishes that the
framework operates at multiple levels of structural realism with
different content at each level. The hierarchy is two-dimensional: an
observation-hierarchy axis (vertical depth --- from foundational to
specific) and a gauge-hierarchy axis (horizontal breadth --- from narrow
to wide gauge equivalence). The combination produces a structure of
realisms operating at multiple intersections.

At the broadest level --- universality-class equivalence across
structural classes (Level G4 in {[}Structure §2{]}) operating at the
partial-trace observational features sub-class (Level C in {[}Structure
§2{]}) --- what is real is the algebra-channel structure of
partial-trace observation: the Born rule, channel-level unitarity,
non-Markovian marginal, and commutant gauge-invariance pattern that any
embedded-observer system must respect. These features are universal
across observer-admitting substrata.

At intermediate levels, realism applies to the SM gauge group (forced
uniquely in OI but shared with any SM-reproducing universality class)
and the algebra-channel structure within the SM-reproducing sub-class.

At the most class-specific level --- the equivalence class
\([(S, \varphi)] / \mathcal{G}_{\text{sub}}\) within OI's structural
class (Level G3 in {[}Structure §2{]}) operating at OI's specific
universality-class representative (Level D in {[}Structure §2{]}) ---
what is real is what this paper articulates: the substratum modulo
\(\mathcal{G}_{\text{sub}}\). OI's specific predictions (Cabibbo angle
\(1/(\pi\sqrt{2})\), Koide ratio \(2/3\), dark-sector phenomenology) are
real at this level, distinguishing OI's universality-class
representative from alternatives within the same partial-trace-features
sub-class.

These are not competing realisms; they are realism at different levels
of the hierarchy. The framework commits to all of them simultaneously,
with content at each level being class-specific (more specific at deeper
levels) or class-universal (more universal at higher levels). What is
real at the universality-class level is what observation extracts from
any partial-trace operation; what is real at the substratum-class level
is the specific equivalence class that the reconstruction theorem
identifies.

This refinement does not weaken the substratum-class realism articulated
above; it situates it within a broader hierarchical structure. The
framework's empirical content (twenty-two quantitative predictions,
specific cubic-lattice substratum, specific gauge group derivation) is
concentrated at the deepest level of the hierarchy --- Level D × Level
G3 --- which is where the reconstruction theorem and substratum gauge
group operate. The broader levels of the hierarchy are where the
framework connects to other unification programs and where
universality-class structural realism applies. The full articulation of
the hierarchical structure is in {[}Structure §2{]}; the
substratum-class realism developed in this paper is one specific level
within that broader structure.

\hypertarget{the-ontological-hierarchy}{%
\subsubsection{6.2 The ontological
hierarchy}\label{the-ontological-hierarchy}}

The triple (S, φ, V) generates every concept in fundamental physics, not
as independent substances but as different aspects of the same
structure. \textbf{Space} is the coupling structure of φ --- the graph
G\_φ determined by which degrees of freedom affect which others ({[}SM,
§2.4{]}). \textbf{Matter} is the state --- localized patterns that
propagate through the coupling graph. \textbf{Energy} is the rate of
change under iteration. \textbf{Time} is the iteration itself.
\textbf{Quantum mechanics} is the observer's compressed description of
the visible sector. \textbf{General relativity} is the thermodynamic
limit of the coupling structure. \textbf{Conservation laws} are
emergent: energy conservation (Noether) is what information conservation
(bijectivity) looks like in the emergent quantum description. None of
these are independent entities; they are descriptions of (S, φ, V) at
different scales.

\hypertarget{the-measurement-problem}{%
\subsubsection{6.3 The measurement
problem}\label{the-measurement-problem}}

On the structural reading, the measurement problem is dissolved. The
wave function is not a component of (S, φ, V) --- it is a derived
object. Since it is derived, not fundamental, asking ``does it
collapse?'' is asking about the behavior of a compression artifact. In
the double-slit experiment, the particle traverses a single slit in the
deterministic substratum. In Wigner's friend, the Friend has a definite
outcome; Wigner's superposition reflects his epistemic deficit.

Branching is forbidden by the rigidity of φ. A fixed bijection on a
finite set has exactly one trajectory from any initial state. There is
no point at which the trajectory splits. The appearance of branching in
the emergent quantum description reflects the observer's uncertainty
about which trajectory they are on (because they cannot see the hidden
sector), not a physical splitting of worlds.

Non-locality in Bell correlations is explained by the coupling graph
G\_φ. Two visible sites prepared in a joint state (entangled) have
correlated hidden-sector configurations --- a consequence of the joint
P-indivisible dynamics at preparation {[}Main, §3.3{]}. The correlations
are mediated by the coupling graph, not by any superluminal influence.
The graph structure ensures that the correlations respect the Tsirelson
bound and violate Bell inequalities without violating parameter
independence.

\hypertarget{pre-registered-falsification-conditions}{%
\subsubsection{6.4 Pre-registered falsification
conditions}\label{pre-registered-falsification-conditions}}

The framework's structural content --- the equivalence class
\([(S, \varphi)]/\mathcal{G}_{\text{sub}}\), the emergent Standard Model
structure, and the co-derivation of quantum mechanics and general
relativity --- makes specific commitments whose future observational
status can now be stated explicitly. This subsection pre-registers the
conditions under which the framework would be falsified, separated from
properties of the specific bijection \(\varphi\) that the framework does
not claim to determine.

The distinction matters for two reasons. First, an explicitly
pre-registered falsification schedule is the empirical counterpart of
the methodological discipline developed in §§3--5: a framework that
derives rather than posits must be willing to specify what would
invalidate the derivation. Second, pre-registration separates structural
commitments --- which the framework stands or falls on --- from
solution-specific properties that lie outside its scope, and which
continued openness on does not undermine the framework.

\textbf{Class A: Structural theorems.} The following are claims about
the equivalence class itself and would falsify the framework's
reconstruction if violated.

\begin{itemize}
\tightlist
\item
  \emph{Exact \(\bar\theta = 0\) at all energy scales.} The strong-CP
  phase is zero at the substratum level (SM §5.3, Theorem 21); any
  observation of a non-zero neutron electric dipole moment consistent
  with \(\bar\theta \gtrsim 10^{-11}\), at sufficient confidence to
  exclude experimental systematics, falsifies the framework's
  T-symmetric substrate commitment.
\item
  \emph{Gauge group SU(3) × SU(2) × U(1) exactly.} Any observation of a
  stable beyond-Standard-Model gauge structure --- an additional \(Z'\)
  interpretable as a new gauge boson rather than a composite resonance,
  stable exotic matter transforming under a novel gauge group, or any
  confirmed extension of the visible-sector gauge content --- falsifies
  the cubic-group decomposition of SM §4.6.
\item
  \emph{Three fermion generations.} Any observation of a stable
  fourth-generation fermion with Standard Model quantum numbers
  falsifies the taste-reduction derivation (SM §4.7). Non-stable exotic
  states consistent with composite or resonance interpretation do not.
\item
  \emph{Majorana neutrinos with normal ordering.} Inverted mass
  ordering, confirmed to high significance, falsifies the taste-breaking
  derivation of the neutrino mass matrix (SM §8.4). Confirmed Dirac
  nature of the neutrino (e.g., a null result for neutrinoless
  double-beta decay at the sensitivity that would require Majorana
  masses well below the framework's prediction) falsifies the prediction
  that no right-handed neutrino exists in the spectrum.
\item
  \emph{Tsirelson bound.} Any observation of a CHSH parameter exceeding
  \(2\sqrt{2}\) in a loophole-free Bell test falsifies the
  P-indivisibility structure ({[}Main, §3.3{]}) --- the bound is a
  theorem of the framework, not an input.
\item
  \emph{Leading Lorentz-invariance-violation coefficient.} The
  cubic-lattice dispersion relation ({[}SM, §4.4{]} Theorem 4,
  generalized to d = 3:
  \(\cos(\omega\epsilon) = \frac{1}{3}\sum_{i=1}^3 \cos(k_i \epsilon)\)
  for a massless mode) yields standard emergent dispersion
  \(\omega = |k|\) at leading order. Series expansion gives the first
  correction as a direction-dependent quadratic term
  \(\omega^2 = k^2[1 - 2\delta(\hat k)(k\epsilon)^2 + O((k\epsilon)^4)]\),
  with \(\delta(\hat k) = \frac{1}{24} - \frac{1}{72}\sum_i \hat k_i^4\)
  evaluated on the rigid \(\mathbb{Z}^3\) representative. The
  rigid-lattice anisotropy (values \(\delta_{[100]} = 1/36\),
  \(\delta_{[111]} = 1/27\)) is \emph{not} a gauge-invariant prediction
  --- \(\mathcal{G}_{\text{sub}}\) generator (iv) of §4 identifies
  \(\mathbb{Z}^3\) with any bounded-degree random graph of the same
  polynomial growth exponent and statistical isotropy (E4), which has no
  preferred axes. The \emph{orientation-averaged} coefficient
  \(\langle\delta\rangle = 1/30\) is gauge-invariant; equivalently, the
  framework predicts the emergent dispersion
  \[\omega^2 = k^2\left[1 - \frac{2}{15}\left(\frac{E}{M_{\text{Pl}}}\right)^2 + O\left(\left(\frac{E}{M_{\text{Pl}}}\right)^4\right)\right]\]
  after using \(\epsilon = 2 l_p\) and
  \(k\epsilon = 2 E/M_{\text{Pl}}\). Three structural features: (i)
  \emph{subluminal}, since \(\delta(\hat k) > 0\) everywhere on the
  sphere (\(\langle\delta\rangle > 0\) is invariant under
  orientation-averaging); (ii) \emph{no linear-in-\(E/M_{\text{Pl}}\)
  term}, because the even-parity cosines in the dispersion admit no
  \(\mathcal{O}(k^3)\) contribution --- any demonstration of linear LIV
  at any scale therefore falsifies the framework's substrate-level
  commitment to the cubic-lattice wave equation; (iii) \emph{specific
  coefficient} \(2/15\), fixed by the cubic-lattice coordination and
  isotropy-gauge average, not tunable. Current quadratic-LIV bounds from
  GRB time-of-flight reach only \(E_{\text{QG}} \gtrsim 10^6\) GeV
  (roughly thirteen orders of magnitude below \(M_{\text{Pl}}\)), so the
  \(2/15\) coefficient is a structural prediction not yet near-term
  testable at the magnitude level; linear-LIV tests, by contrast, are
  already at sensitivities where any positive detection falsifies.
\end{itemize}

\textbf{Class B: Parameter-free retrodictions.} The framework produces a
set of numerical predictions that follow from the structural content
with no adjustable parameters beyond a small set of acknowledged
empirical inputs (\(m_s\) for the overall fermion mass scale, \(m_t\)
for the Higgs RGE running, \(\eta\) for the Jarlskog phase,
\(Q_{\text{down}}\) for the down-sector Koide normalization). These
predictions are stratified into four sub-classes reflecting their
structural status:

\emph{Class B-S (unconditional structural).} Six predictions are Layer 0
or Layer 1 unconditional structural retrodictions:

\begin{itemize}
\tightlist
\item
  The Cabibbo angle \(\lambda = 1/(\pi\sqrt{2})\) ({[}SM §7.1{]}) ---
  the chirality verification of the taste-changing vertex's spinor
  structure is closed via Mason et al.~(HPQCD, hep-lat/0209152),
  establishing that taste-changing transitions in staggered quarks have
  spinor structure given by a combination of \(\gamma_\mu\) and
  \(\gamma_{5\mu}\) (both chirality-preserving).
\item
  The Wolfenstein parameter \(A = \sqrt{2/3}\) ({[}SM §7.1{]}).
\item
  The mass ratio \(m_d/m_s = 1/(2\pi^2)\) via the Gatto-Sartori-Tonin
  relation applied to the Layer 1 Cabibbo derivation ({[}SM §7.1{]}).
\item
  The Koide angle \(\theta_0 = C_2/d^2 = 2/9\) ({[}SM §7.2{]}) --- the
  sharpest empirical match in the framework (0.02\%) and the cleanest
  structural derivation (cubic-group quadratic Casimir over bandwidth
  squared, with the §7.2 uniqueness table enumerating ten alternative
  dimensionless ratios that fail to match).
\item
  The Bekenstein-Hawking coefficient \(S = A/(4 l_p^2)\) with the factor
  \(1/4\) derived as the ratio \(2\pi/(8\pi)\) between the Euclidean KMS
  period and the Jacobson Einstein-equation prefactor ({[}GR §5{]}).
  Confirmed at 99.999\% by the GW250114 area-law observation.
\item
  The MOND critical acceleration \(a_0 = cH/6\) ({[}GR §7.3{]}) --- the
  \(1/6\) geometric factor is closed via Verlinde 2016 (eq. 1.7) as
  \((d-3)/[(d-2)(d-1)] = 1/6\) in \(d = 4\), with OI providing the
  structural grounding (volume-law entropy mechanism derived from
  condition C2 rather than postulated). Dependent predictions (baryonic
  Tully-Fisher, crossover radius \(r_M\)) inherit B-S status.
\end{itemize}

Any of these moving outside \(3\sigma\) on precision upgrade falsifies
the framework directly with no recoverable error mode.

\emph{Class B-L (layered conditional, pending open derivations).} Three
predictions follow from the structural form combined with one Layer-2
substrate input or empirical input. Each has an explicit derivational
gap whose closure would move the prediction to B-S:

\begin{itemize}
\tightlist
\item
  The mass ratio \(m_u/m_d = \sqrt{\theta_0} = \sqrt{2/9}\) ({[}SM
  §7.2{]}) pending the explicit derivation of the ``different channels''
  mechanism producing the square root, structurally analogous to the
  open OI-vertex 1-loop derivation for \(K = 1/2\) ({[}SM §7.5{]}) in
  being mixed-layer --- substratum cubic-group structure combined with
  electroweak-symmetry-breaking emergent machinery.
\item
  The mass ratio \(m_b/m_\tau = 4.28/Z_S\) ({[}SM §7.5{]}) pending the
  explicit OI-vertex 1-loop computation confirming \(K = 1/2\), together
  with four bridge gaps in the SχPT inheritance ({[}SM §7.5{]}): the
  substratum measure reduction under Theorem 2, the 3D MC vs 4D SχPT
  dimensional bridge, the 3D BZ vs 4D Dirac taste-structure equivalence,
  and the 3D vs 4D parity identification of chirality. The per-vertex
  \(\cos^2\) symmetrization producing \(K = 1/2\) is
  dimension-independent and structurally robust; the four bridges are
  derivation gaps inherited by the other §7 mass-cluster predictions
  that share the SχPT machinery. The closure path is the two-loop SχPT
  calculation specializing Panagopoulos--Spanoudes 2017 to the OI
  cubic-group setting.
\item
  The PMNS predictions \(\sin^2\theta_{12} = 1/3 - 1/(4\pi^2)\) and
  \(\sin^2\theta_{23} = 1/2 + 1/(2\pi^2)\) ({[}SM §7.3{]}) pending
  derivation of Cond 2 (the structural relation among \(A_2\) Wilson
  coefficients) from cubic-lattice Yukawa structure. The reactor angle
  \(\sin^2\theta_{13}\) is independent of Cond 2 and is in Class B-S.
\end{itemize}

Class B-L violations at \(3\sigma\) on precision upgrade are addressed
under the pre-registration commitment via three branches: (a) closure of
the open derivation \emph{confirms} the prediction (recoverable from a
transient data fluctuation); (b) closure \emph{fails to confirm} the
prediction while preserving the framework's structural commitments
(recoverable error in the specific derivation, with the structural
commitments intact); (c) demonstration that no consistent OI-internal
derivation can produce the observed value (fatal, framework-falsifying).

\emph{Class B-M (mass chain).} Two predictions inherit a single
empirical input scale through a structural relation:

\begin{itemize}
\tightlist
\item
  The electron and muon masses \(m_e\) and \(m_\mu\) via the Koide chain
  from the empirical input \(m_\tau\) ({[}SM §7.2{]}). The prediction is
  structurally \(\theta_0 = 2/9\) with \(Q = 2/3\); the absolute scale
  is set by \(m_\tau\).
\item
  The Higgs mass \(m_H\) via SM RGE running from the structural boundary
  condition \(\lambda(M_{\text{Pl}}) = 0\) ({[}SM §7.4{]}) with \(m_t\)
  as the empirical input.
\end{itemize}

Class B-M predictions falsify under the same precision-upgrade protocol
as B-L, but the inheritance is from an empirical input scale rather than
from an open derivation.

\emph{Class B-R (retrodictions).} Two gauge-coupling values are
explicitly retrodictions per the §6.3 {[}SM{]} parameter count:

\begin{itemize}
\tightlist
\item
  The SU(2) gauge coupling at \(M_Z\) ({[}SM §6, items 11-12 in
  §7.6{]}).
\item
  The SU(3) gauge coupling at \(M_Z\).
\end{itemize}

Three fitted parameters \((\delta_0, A, B)\) against three observed
couplings produces zero residual by construction; these classifications
are honest. Class B-R predictions can never falsify in the
precision-upgrade sense --- their value is what falsification of the
\emph{framework's gauge-sector structural commitments} (universality of
\(\delta_0\), the resummation form, the structural prediction
\(1/\alpha_0 = 23.25\) from the 1-loop staggered VP, and the
\(A \cdot B = 48.0 \pm 1.5\) cross-check) would test.

\emph{Status summary.} Of the predictions in Class B: \textbf{4 are
unconditional structural} (B-S), \textbf{6 are layered conditional
pending open derivations} (B-L), \textbf{2 are mass-chain inheritances}
(B-M), and \textbf{2 are retrodictions} (B-R). Total 14 predictions in
Class B. The empirical match holds for all 14 within \textasciitilde1\%
or \textasciitilde1σ.

The current status of each is documented in the SM and GR companion
papers with experimental references. As experimental precision improves,
the discrimination power of the B-S set increases monotonically; closure
of the open derivations would move B-L predictions to B-S over time.

\textbf{Class C: Framework-level commitments.} Two broader commitments
are implicit in the construction and would falsify it if contradicted.

\begin{itemize}
\tightlist
\item
  \emph{Classical-memory simulability of any genuinely quantum process.}
  The characterization theorem ({[}Main, §3.4{]}) establishes that
  quantum mechanics is the description produced by an embedded observer
  satisfying C1--C3 on a deterministic substrate with bounded classical
  hidden-sector memory. This implies every process the framework calls
  quantum admits a classical-memory simulation at appropriate scale; any
  demonstration of a fundamental quantum process that provably cannot be
  simulated by finite classical memory under any relabeling falsifies
  the framework at the characterization-theorem level. The commitment is
  specific: non-Markovian models are permitted and expected (they are
  how the hidden-sector correlation persistence manifests), but they
  must be realizable by finite classical memory. A process whose
  non-Markovianity is provably supra-classical falsifies.
\item
  \emph{The dark sector as entropy-displacement, not as a novel particle
  content.} The framework's account of dark energy as a Type II RVM
  functional form and dark matter as frozen boundary entropy (GR §7.3)
  is incompatible with the direct detection of a WIMP interpretable as a
  fundamental particle with Standard-Model-like interactions and no
  entropy-displacement signature. A confirmed WIMP with such properties
  falsifies the framework's dark-sector account while leaving the rest
  of the structure intact --- it would require abandoning the §7.3
  derivation and either identifying an error in it or accepting that the
  framework's dark-sector story is wrong even as other derivations
  survive.
\end{itemize}

\textbf{Cosmological commitment.} The dark energy sector is subtler than
Classes A--C. The framework commits to the Type II running-vacuum
functional form (GR §7.1), with magnitude given at theorem level by the
emergent-QFT calculation
\(|\nu_{\rm QFT}| \sim (M_{\rm SM}/M_{\rm Pl})^2 \sim 10^{-32}\),
observationally indistinguishable from zero. Substratum contributions
beyond this are constrained by the boundary mode count of §3-§5, which
catalogues modes by spatial location alone (one mode per \(\epsilon^2\)
cell, no inherent frequency assignment); the natural 2D distribution
concentrates them near the UV cutoff, contributing well below
\(|\nu_{\rm QFT}|\). The framework's theorem-level magnitude prediction
is therefore \(|\nu| \approx |\nu_{\rm QFT}|\), with the Bertini et
al.~(2025) measurement \(\nu = -(2.5 \pm 1.3) \times 10^{-4}\)
consistent with this prediction at \(2\sigma\), where Bertini is itself
consistent with zero. The falsification band is correspondingly
stratified: a DESI Year 5 measurement of \(|\nu| \gtrsim 10^{-4}\) at
high significance would require additional substratum structure
(log-uniform mode-per-decade distribution at the cosmological horizon)
not currently in §3-§5; a measurement of \(\nu\) with high-significance
positive sign would falsify the Type II functional form itself, which is
a structural commitment rather than a magnitude one. The first outcome
is a refinement within the current bracket; the second is a structural
falsification.

\textbf{What does not falsify.} The framework makes no commitment to the
specific values of solution-specific quantities, and their future
measurement does not bear on the framework's structural status. These
include: the absolute scale of fermion masses (the choice of \(m_s\) as
empirical input is acknowledged in SM §7.6); the CP-violating phases in
CKM and PMNS (solution-specific properties set by the basis mismatch
between up-type and down-type Yukawa mass eigenstates, compatible with
the framework's \(\bar\theta = 0\) theorem per {[}SM, §5.4{]}); the
numerical value of the Type II \(\nu\) coefficient within the bracketed
range above; initial conditions at any cosmic time (the framework
derives the structure from which evolution proceeds, not the state at
\(t_0\)); and the specific microstate \(\varphi\) within the equivalence
class (Theorem 24 classifies this as gauge, so the question of which
representative our universe is in is provably without empirical
content). A framework is not falsified by the presence of
solution-specific inputs; it is falsified by failure of its structural
commitments. The Class A--C list, together with the stratified
dark-energy band, is the complete structural commitment of the framework
as currently written.

\textbf{The pre-registration commitment.} The author commits to treating
the Class A conditions, the Class C commitments, and a Class B
retrodiction exceeding \(3\sigma\) on precision upgrade as
framework-falsifying rather than as an invitation to post-hoc rescue. In
the case of a Class B violation, the framework will be revisited to
determine whether the violation reflects a computational error in the
specific retrodiction, an incorrect specific-bijection input, or a
structural problem with the derivation; the first two are recoverable,
the third is fatal. The framework does not reserve the right to add
structural postulates after the fact to accommodate disconfirming
evidence. Pre-registration of this commitment is itself part of the
falsifiability content of the framework.

\textbf{Regimes of deliberate silence.} Three regimes lie outside the
framework's predictive scope by construction and are not objects of
falsification claims. First, the \emph{pre-horizon early universe} ---
the regime before a stable causal partition with \(\tau_B \sim H^{-1}\)
exists --- has no observer in the framework's sense ({[}Main, §1.3{]})
and no horizon-level structure for the {[}GR{]} derivations to act on;
the framework's only commitment in this regime is a boundary condition
on deeper theories, namely that any successful pre-horizon dynamics must
produce a C1--C3-satisfying regime by the time a horizon stabilizes.
Second, the \emph{equilibrium phase} of \((S, \varphi)\) is structurally
observer-free by the structural observer-selection theorem ({[}Main,
§4.6{]}); the framework makes no prediction for this regime because it
has no observers to whom a prediction could be addressed. Third, regions
\emph{beyond the observer's causal horizon} are part of that observer's
hidden sector and are addressed only through boundary effects (the
area-law entropy of {[}GR, §5{]}); the framework takes no
observer-independent position on their internal content. These silences
are structural rather than methodological --- they follow from the
characterization theorem's identification of QM with embedded
observation, and from the observer-selection theorem's restriction of
observers to the non-equilibrium phase --- and stating them explicitly
sharpens the scope of the positive commitments above.

\begin{center}\rule{0.5\linewidth}{0.5pt}\end{center}

\hypertarget{conclusion}{%
\subsection{7. Conclusion}\label{conclusion}}

The substratum-level results developed in this paper --- the
reconstruction theorem (Theorem 23) and the substratum gauge group
(Theorem 24) --- turn the framework's three derivations into a single
construction. The reconstruction theorem establishes that observed
quantum mechanics with Bell violations, finite boundary entropy, and
spatial isotropy, together with the framework's structural assumptions,
uniquely determine the equivalence class
\([(S, \varphi)]/\mathcal{G}_{\text{sub}}\) at the lattice level, with
the Standard Model gauge group and \(\bar\theta = 0\) as forced
retrodictions, and the anomaly-free hypercharge assignment forced once
the observed family pattern is given (fermion embedding closed via the
link-carrier construction, with the absolute generation count locked at
exactly three by coupling-degree minimality). The substratum gauge group
identifies the kernel of this inverse map and shows that the Standard
Model gauge group is its visible-sector shadow. Together, these two
results establish that the framework's emergence of quantum mechanics
{[}Main{]}, its derivation of the Standard Model {[}SM{]}, and its
derivation of the gravitational sector {[}GR{]} are not three
independent applications of the same trace-out machinery but three
projections of one object --- the bijection \((S, \varphi)\) --- viewed
at different levels of description.

The synthesis claim that follows is not a unification of quantum
mechanics and general relativity in the traditional sense. The
traditional unification problem treats the two theories as competing
accounts of the same regime and asks how to merge them. The present
framework treats them as projections of the same construction onto two
different substructures --- the visible-sector trace-out and the
boundary-thermodynamic limit --- and the question of how to merge them
is dissolved as a category error rather than answered as a technical
problem. The result is not a theory of everything but a structural
account of why the apparent incompatibility between quantum mechanics
and general relativity has been so resistant to resolution: there is
nothing to resolve, because the two are not in competition. They are
co-derived from the same object by the same procedure under the same
conditions, and the framework developed across this four-paper synthesis
({[}Main{]}, {[}SM{]}, {[}GR{]}, {[}Substratum{]}) makes this exact
rather than approximate.

Three sets of open problems remain. Neither these results nor this
four-paper synthesis require the framework to be the final theory of
physics: the substratum may itself be an effective description of
something deeper, and the trace-out may have corrections beyond the
leading-order results developed here. First, the \emph{specific
bijection} \(\varphi\) that describes our universe is not determined by
the framework --- only the equivalence class
\([(S, \varphi)]/\mathcal{G}_{\text{sub}}\) is. The framework predicts
the structural features of the Standard Model and the gravitational
sector but not the specific values of (for example) the lightest fermion
mass or the specific numerical value of the Type II RVM coefficient
\(\nu\) beyond its dimensional order of magnitude and sign. These are
properties of the particular bijection, analogous to the mass of the sun
in general relativity. Second, the \emph{trans-Planckian regime} --- the
regime in which the lattice spacing \(\epsilon = 2 l_p\) is not small
compared to the relevant length scale --- lies outside the framework's
leading-order results. The framework presents the lattice as
fundamental, not as a regulator approximating a continuum theory, and so
does not need a continuum limit in the standard sense; but the
corrections to the leading-order results in the regime where the
partition geometry varies on lattice scales are an open question. Third,
the \emph{initial conditions} --- the specific configuration of the
bijection at any given time --- are likewise not determined by the
framework. The framework establishes which structures emerge from the
bijection but not which microstate the universe is currently in.

These open problems are framed within the construction rather than
against it. Each is sharply formulated and admits a definite (if
presently unanswered) form. The progress reported here is that the
structural relationship between quantum mechanics and general relativity
--- the central open problem of fundamental physics for nearly a century
--- is fixed by the construction developed across this four-paper
synthesis, and the remaining open problems are problems within the
construction rather than problems with it.

\begin{center}\rule{0.5\linewidth}{0.5pt}\end{center}

\hypertarget{references}{%
\subsection{References}\label{references}}

{[}1{]} R. Bousso, ``The holographic principle,'' \emph{Rev.~Mod. Phys.}
\textbf{74}, 825 (2002).

{[}2{]} N. Bao, S. M. Carroll, and A. Singh, ``The Hilbert space of
quantum gravity is locally finite-dimensional,'' \emph{Int. J. Mod.
Phys. D} \textbf{26}, 1743013 (2017).

{[}3{]} J. Maldacena and L. Susskind, ``Cool horizons for entangled
black holes,'' \emph{Fortsch. Phys.} \textbf{61}, 781--811 (2013);
arXiv:1306.0533.

\begin{center}\rule{0.5\linewidth}{0.5pt}\end{center}

\textbf{Companion papers (cited inline by short name):}

{[}Main{]} A. Maybaum, ``The Incompleteness of Observation,'' (2026).

{[}SM{]} A. Maybaum, ``The Standard Model from a Cubic Lattice,''
(2026).

{[}GR{]} A. Maybaum, ``ℏ, the Bekenstein-Hawking Entropy, and Dynamical
Dark Energy from the Cosmological Horizon,'' (2026).

\end{document}
